% ======================================================================
% 第 3 章:狄拉克场
% (原书第 35-76 页;PDF 第 58-99 页)
% ======================================================================
\chapter{狄拉克场}

在详尽处理了最简单的相对论性场方程之后,我们现在转向第二简单的方程——狄拉克方程。你可能已经熟悉狄拉克方程最初的形式,即作为单粒子量子力学波动方程。%
\footnote{这一主题参见,例如,Schiff (1968) 第~13 章;Baym (1969) 第~23 章;Sakurai (1967) 第~3 章。虽然本章自成一体,但我们建议你在某个时候也学习一下单粒子狄拉克方程。} 在本章中,我们的观点将完全不同。首先,我们将把狄拉克方程重新推导为经典的相对论性场方程,特别强调其相对论不变性。然后,在第 3.5 节,我们将以类似于克莱因-戈登场所用的方式来量子化狄拉克场。

\section{波动方程中的洛伦兹不变性}

首先我们必须处理一个在第 2 章中一笔带过的问题:当我们说一个方程``相对论不变''时,是什么意思?一个合理的定义如下:如果 $\phi$ 是一个场或场集合,$D$ 是某个微分算符,那么陈述``$D\phi = 0$ 是相对论不变的''意味着:如果 $\phi(x)$ 满足这个方程,并且我们旋转或洛伦兹加速到不同的参考系,那么变换后的场在新参考系中满足相同的方程。等价地,我们可以想象在物理上以某个公共角度或速度旋转或洛伦兹加速所有粒子或场;同样,方程 $D\phi = 0$ 在变换之后应当成立。在下面的分析中,我们将对变换采用这种``主动''的观点。

场论的拉格朗日表述使得讨论洛伦兹不变性特别容易。按照上述定义,如果运动方程来源于一个洛伦兹标量的拉格朗日量,那么它自动是洛伦兹不变的。这是最小作用量原理的直接结果:如果洛伦兹加速使拉格朗日量保持不变,那么作用量极值的洛伦兹加速将是另一个极值。

作为一个例子,考虑克莱因-戈登理论。我们可以把任意洛伦兹变换写为
\begin{equation}
  x^\mu \to x'^{\mu} = \Lambda^\mu_{\ \nu} x^\nu,
  \label{eq:3.1}
\end{equation}
其中 $\Lambda$ 是某个 $4\times4$ 矩阵。克莱因-戈登场 $\phi(x)$ 在这个变换下会怎样?把场 $\phi$ 想成是在测量某个分布于空间中的量的局域值。如果这个量在 $x=x_0$ 处有聚集,$\phi(x)$ 将在 $x_0$ 处有极大值。如果我们现在对原始分布进行洛伦兹加速,新的分布将在 $x=\Lambda x_0$ 处有极大值。如图 \ref{fig:3.1}(a) 所示。场相应的变换为
\begin{equation}
  \phi(x) \to \phi'(x) = \phi(\Lambda^{-1}x).
  \label{eq:3.2}
\end{equation}
也就是说,变换后的场在经过洛伦兹加速的点上取值,等于原始场在加速前的点上取值。

\begin{figure}[t]
  \centering
  \includegraphics[width=0.7\textwidth]{fig3.1.png}
  \caption{对矢量场进行旋转时,既会影响矢量的取向,也会影响包含该构型的区域的位置。}
  \label{fig:3.1}
\end{figure}

我们应该检验这个变换是否使克莱因-戈登拉格朗日量的形式保持不变。根据 \eqref{eq:3.2},质量项 $\frac{1}{2}m^2\phi^2(x)$ 只是被平移到点 $(\Lambda^{-1}x)$。$\partial_\mu\phi(x)$ 的变换为
\begin{equation}
  \partial_\mu\phi(x) \to
    \frac{\partial}{\partial x^\mu}\bigl(\phi(\Lambda^{-1}x)\bigr)
    = (\Lambda^{-1})^{\nu}_{\ \mu}(\partial_\nu\phi)(\Lambda^{-1}x).
  \label{eq:3.3}
\end{equation}
由于度规张量 $g_{\mu\nu}$ 是洛伦兹不变的,矩阵 $\Lambda^{-1}$ 满足恒等式
\begin{equation}
  (\Lambda^{-1})^{\mu}_{\ \rho}(\Lambda^{-1})^{\nu}_{\ \sigma}
    g_{\mu\nu} = g_{\rho\sigma} .
  \label{eq:3.4}
\end{equation}
利用这个关系,我们可以计算克莱因-戈登拉格朗日量动能项的变换律:
\begin{align}
  (\partial_\mu\phi(x))^2 &\to
    g^{\mu\nu}(\partial_\mu\phi'(x))(\partial_\nu\phi'(x))
  = g^{\mu\nu}\bigl[(\Lambda^{-1})^{\rho}_{\ \mu}\bigr]
    \bigl[(\Lambda^{-1})^{\sigma}_{\ \nu}\bigr]
    (\partial_\rho\phi)(\partial_\sigma\phi)(\Lambda^{-1}x)\\
  &= g^{\rho\sigma}(\partial_\rho\phi)(\partial_\sigma\phi)(\Lambda^{-1}x).
\end{align}
因此,整个拉格朗日量只是作为一个标量变换:
\begin{equation}
  \mathcal{L}(x) \to \mathcal{L}(\Lambda^{-1}x).
  \label{eq:3.5}
\end{equation}
由 $\mathcal{L}$ 对时空积分构成的作用量 $S$ 是洛伦兹不变的。类似的计算表明运动方程是不变的:
\begin{equation}
  (\partial^2 + m^2)\phi'(x) = \bigl[
    (\Lambda^{-1})^{\mu}_{\ \lambda}\partial_\lambda
    (\Lambda^{-1})^{\nu}_{\ \sigma}\partial_\sigma + m^2 \bigr]
    \phi(\Lambda^{-1}x) = (g^{\lambda\sigma}\partial_\lambda\partial_\sigma
    + m^2)\phi(\Lambda^{-1}x) = 0 .
\end{equation}
用于 $\phi$ 的变换律 \eqref{eq:3.2} 是一个场最简单的变换律。它只对单分量场是唯一可能的形式。但我们知道有多分量场的例子,它们以更复杂的方式变换。最熟悉的例子是矢量场,例如 4-流密度 $j^\mu(x)$ 或矢量势 $A^\mu(x)$。在这种情况下,分布于时空中的量还带有取向,这个取向必须被旋转或洛伦兹加速。如图 \ref{fig:3.1}(b) 所示,当场求值点改变时,取向必须向前旋转:
\begin{equation}
  V^\mu(x) \to (\Lambda^{-1})^{\mu}_{\ \nu}V^\nu(\Lambda^{-1}x)
  \qquad\text{在洛伦兹变换下。}
\end{equation}
任意秩的张量可以通过添加更多指标由矢量构造,变换律中相应地有更多个 $\Lambda$ 因子。利用这样的矢量和张量场,我们可以写出各种各样的洛伦兹不变方程,例如麦克斯韦方程组,
\begin{equation}
  \partial_\mu F^{\mu\nu} = 0
  \qquad\text{或}\qquad
  \partial^2 A^\mu - \partial_\nu\partial^\mu A^\nu = 0,
  \label{eq:3.6}
\end{equation}
它们由拉格朗日量
\begin{equation}
  \mathcal{L}_{\mathrm{Maxwell}} =
    -\tfrac{1}{4}(\partial_\mu A_\nu - \partial_\nu A_\mu)^2 .
  \label{eq:3.7}
\end{equation}
导出。一般地说,任何每一项都含有同一组未缩并洛伦兹指标的方程,在洛伦兹变换下自然是洛伦兹不变的。这种张量记号方法给出了一大类洛伦兹不变方程,但事实证明还有更多。我们如何找到它们?我们可以尝试系统地找出一个场的所有可能的变换律。那样写出不变的拉格朗日量就不难了。为简单起见,我们将把注意力限制在线性变换上,这样,如果 $\Phi_a$ 是一个 $n$ 分量多重态,洛伦兹变换律就由一个 $n\times n$ 矩阵 $M(\Lambda)$ 给出:
\begin{equation}
  \Phi_a(x) \to M_{ab}(\Lambda)\Phi_b(\Lambda^{-1}x).
  \label{eq:3.8}
\end{equation}
可以证明,最一般的非线性变换律可以由这些线性变换构造,因此考虑比 \eqref{eq:3.8} 更一般的变换没有任何好处。在下面的讨论中,我们将略去场自变量的变化,把变换 \eqref{eq:3.8} 写成形式
\begin{equation}
  \Phi \to M(\Lambda)\Phi .
  \label{eq:3.9}
\end{equation}
矩阵 $M(\Lambda)$ 允许的形式有哪些?对 $M(\Lambda)$ 的基本限制是通过设想两个相继的变换 $\Lambda$ 和 $\Lambda'$ 找到的。最终结果必须是一个新的洛伦兹变换 $\Lambda''$;也就是说,洛伦兹变换构成一个群。这给出了矩阵 $M(\Lambda)$ 必须满足的自洽条件:在两个变换的序列下,
\begin{equation}
  \Phi \to M(\Lambda')M(\Lambda)\Phi = M(\Lambda'')\Phi,
  \label{eq:3.10}
\end{equation}
其中 $\Lambda'' = \Lambda'\Lambda$。因此矩阵 $M$ 与变换 $\Lambda$ 之间的对应必须在乘法下保持。用数学语言说,矩阵 $M$ 必须构成洛伦兹群的一个 $n$ 维表示。所以我们现在的问题用数学语言表述为:洛伦兹群的(有限维)矩阵表示有哪些?

在回答洛伦兹群的这个问题之前,让我们考虑一个更简单的群:三维旋转群。这个群具有每一个维度 $n$ 的表示,在量子力学中它们是旋转不同自旋粒子的 $n$ 分量波函数的矩阵,因而是熟悉的。维度通过 $n = 2s+1$ 与自旋量子数 $s$ 相联系。最重要的非平凡表示是二维表示,对应自旋 $1/2$。这个表示的矩阵是行列式为 1 的 $2\times2$ 幺正矩阵,可以表示为
\begin{equation}
  U = e^{-i\theta^l\sigma^l/2},
  \label{eq:3.11}
\end{equation}
其中 $\theta^l$ 是三个任意参数,$\sigma^l$ 是泡利西格马矩阵。

对任何连续群,位于恒等元无穷小邻域内的变换定义了一个矢量空间,称为该群的李代数。这个矢量空间的基矢量称为李代数或群的生成元。对旋转群,生成元是角动量算符 $J^i$,它们满足对易关系
\begin{equation}
  [J^i, J^j] = i\epsilon^{ijk}J^k .
  \label{eq:3.12}
\end{equation}
有限旋转操作通过对这些算符取指数构成:在量子力学中,算符
\begin{equation}
  R = \exp\Bigl[-\frac{i}{\hbar}\bm{\theta}\cdot\mathbf{J}\Bigr]
  \label{eq:3.13}
\end{equation}
给出绕轴 $\bm{\theta}$ 转过角度 $|\bm{\theta}|$ 的旋转。算符 $J^i$ 的对易关系决定了这些旋转算符的乘法律。因此,满足对易关系 \eqref{eq:3.12} 的一组矩阵,通过如 \eqref{eq:3.13} 的取指数,就产生旋转群的一个表示。在上段给出的例子中,角动量算符的表示
\begin{equation}
  J^l = \frac{\sigma^l}{2}
  \label{eq:3.14}
\end{equation}
产生式 \eqref{eq:3.11} 给出的旋转群的表示。一般来说,我们可以通过寻找群的生成元的矩阵表示(它们必须满足正确的对易关系),然后对这些无穷小变换取指数,来找到连续群的矩阵表示。

对我们当前的问题,我们需要知道洛伦兹变换群的生成元的对易关系。对旋转群,可以把生成元写成微分算符来推出对易关系;从表达式
\begin{equation}
  \mathbf{J} = \mathbf{x}\times\mathbf{p} = \mathbf{x}\times(-i\nabla),
  \label{eq:3.15}
\end{equation}
角动量对易关系 \eqref{eq:3.12} 可以直接得出。在 \eqref{eq:3.15} 中使用叉积是三维情形的特殊之处。不过,我们也可以把算符写成一个反对称张量,
\begin{equation}
  J^i = -\tfrac{1}{2}\epsilon^{ijk}J^{jk},
\end{equation}
使得 $J^3 = J^{12}$ 等等。现在,对四维洛伦兹变换的推广就很自然了:
\begin{equation}
  J^{\mu\nu} = i(x^\mu\partial^\nu - x^\nu\partial^\mu).
  \label{eq:3.16}
\end{equation}
我们很快将看到,这六个算符生成洛伦兹群的三个洛伦兹加速和三个旋转。

为了确定洛伦兹代数的对易规则,我们现在只需计算微分算符 \eqref{eq:3.16} 的对易子。结果是
\begin{equation}
  [J^{\mu\nu}, J^{\rho\sigma}] = i\bigl( g^{\nu\rho}J^{\mu\sigma}
    - g^{\mu\rho}J^{\nu\sigma} - g^{\nu\sigma}J^{\mu\rho}
    + g^{\mu\sigma}J^{\nu\rho} \bigr).
  \label{eq:3.17}
\end{equation}
任何表示这个代数的矩阵都必须满足同样的对易规则。

为了确认我们弄对了,让我们看一个特殊的表示(我们将简单地把它从帽子里变出来)。考虑 $4\times4$ 矩阵
\begin{equation}
  (\mathcal{J}^{\mu\nu})^{\alpha}_{\ \beta} =
    i\bigl( \delta^{\mu\alpha}\delta^{\nu}_{\ \beta}
      - \delta^{\nu\alpha}\delta^{\mu}_{\ \beta} \bigr).
  \label{eq:3.18}
\end{equation}
不难验证这些矩阵满足对易关系 \eqref{eq:3.17}。我们可以把它们想成对 4-矢量的一组线性变换,以我们已经看到的方式旋转或洛伦兹加速 4-矢量 $V^\mu$。于是,矢量场上的有限洛伦兹变换为
\begin{equation}
  V^\mu \to \Bigl[ \exp\bigl(
    -\tfrac{i}{2}\Omega_{\rho\sigma}\mathcal{J}^{\rho\sigma} \bigr)
    \Bigr]^{\mu}_{\ \nu} V^\nu ,
  \label{eq:3.19}
\end{equation}
其中 $\Omega_{\rho\sigma}$ 是变换的参数,选择得使 $\Lambda = \exp(-\frac{i}{2}\Omega_{\rho\sigma}\mathcal{J}^{\rho\sigma})$。

现在我们需要旋量的对应表示。对自旋 $1/2$,维度 $n$ 是 2,所以我们寻找满足相同对易关系 \eqref{eq:3.17} 的 $2\times2$ 矩阵,其中生成元表示三维旋转的角动量 $J^i = \sigma^i/2$。我们现在可以看到构造完整洛伦兹代数生成元的系统方法:把六个生成元用泡利矩阵写出,然后寻找一个 $2\times2$ 表示。三个旋转生成元就是自旋算符
\begin{equation}
  S^{ij} = \frac{1}{2}\epsilon^{ijk}\sigma^k
  \label{eq:3.20}
\end{equation}
(比较式 \eqref{eq:3.14});这重现了旋转子群的交换子 \eqref{eq:3.12}。三个洛伦兹加速生成元必须是 $2\times2$ 矩阵;为了找到它们,我们注意到对易关系
\begin{equation}
  [K^i, K^j] = -i\epsilon^{ijk}J^k
\end{equation}
(由 \eqref{eq:3.17} 推出)必须被满足。可以验证矩阵
\begin{equation}
  K^i = \frac{i}{2}\sigma^i
  \label{eq:3.21}
\end{equation}
满足这个关系。因此,洛伦兹群的一个 $2\times2$ 表示是
\begin{equation}
  \Lambda_{1/2} = \exp\Bigl( -\frac{i}{2}\Omega_{\mu\nu}S^{\mu\nu} \Bigr),
  \label{eq:3.22}
\end{equation}
其中生成元 $S^{\mu\nu}$ 由
\begin{equation}
  S^{\mu\nu} = \begin{cases}
    \frac{1}{2}\epsilon^{ijk}\sigma^k & (\mu,\nu)=(i,j)\\
    \frac{i}{2}\sigma^i & (\mu,\nu)=(0,i)
  \end{cases}
\end{equation}
以泡利矩阵给出。我们注意到 $\sigma^i$ 是厄米的,所以旋转生成元 $S^{ij}$ 是厄米的,而洛伦兹加速生成元 $S^{0i}$ 是反厄米的。因此有限维表示 \eqref{eq:3.22} 不是幺正的(这与非紧致洛伦兹群相称)。不过,这些旋量指标不是我们的主要关注点;重要的是,我们现在有了洛伦兹群的一个候选 $2\times2$ 表示,可以用来写出 2 分量场的变换律。

让我们显式写出按这个表示变换的对象。狄拉克旋量是一个 4 分量对象
\begin{equation}
  \psi(x) = \begin{pmatrix}\psi_1\\ \psi_2\\ \psi_3\\ \psi_4\end{pmatrix},
\end{equation}
我们经常把它拆成两个 2 分量部分,正如下面将要描述的那样。我们不用从头构造洛伦兹代数,而是用另一种但等价的方法。用反对易关系
\begin{equation}
  \{\gamma^\mu, \gamma^\nu\} = 2g^{\mu\nu},
  \label{eq:3.22b}
\end{equation}
定义四个 $4\times4$ 矩阵 $\gamma^\mu$,称为狄拉克代数。(在闵可夫斯基空间中反对易子为 $\{\gamma^\mu,\gamma^\nu\} = \gamma^\mu\gamma^\nu + \gamma^\nu\gamma^\mu$。)可以证明矩阵
\begin{equation}
  S^{\mu\nu} = \frac{i}{4}[\gamma^\mu, \gamma^\nu]
  \label{eq:3.23}
\end{equation}
满足洛伦兹代数的对易关系 \eqref{eq:3.17}。因此 $\gamma$ 矩阵提供了洛伦兹生成元的一个表示,其中狄拉克场 $\psi$ 在洛伦兹变换下的变换律为
\begin{equation}
  \psi(x) \to \Lambda_{1/2}\psi(\Lambda^{-1}x),
  \qquad \Lambda_{1/2} = \exp\Bigl(
    -\frac{i}{2}\Omega_{\mu\nu}S^{\mu\nu}\Bigr).
  \label{eq:3.24}
\end{equation}

对二维闵可夫斯基空间,我们可以找到一组满足狄拉克代数的 $2\times2$ 矩阵:
\begin{equation}
  \gamma^0 = \sigma^3,\qquad \gamma^1 = i\sigma^1,
\end{equation}
使得
\begin{equation}
  \{\gamma^i, \gamma^j\} = -2\delta^{ij}.
  \label{eq:3.24b}
\end{equation}
第一行中的因子 $i$ 和第二行中的负号纯属约定。于是表示洛伦兹代数的矩阵为
\begin{equation}
  S^{ij} = \frac{1}{2}\epsilon^{ijk}\sigma^k
  \label{eq:3.25}
\end{equation}
我们认出它是旋转群的二维表示。

现在让我们来找四维闵可夫斯基空间的狄拉克矩阵 $\gamma^\mu$。事实证明,这些矩阵至少必须是 $4\times4$ 的。(例如,不存在第四个与三个泡利西格马矩阵都反对易的 $2\times2$ 矩阵。)此外,狄拉克代数的所有 $4\times4$ 表示都幺正等价。%
\footnote{这一陈述与前面的陈述可由习题~3.1 中推导的洛伦兹群表示的一般理论得出。} 因此我们只需写出狄拉克代数的一个显式实现。一个 $2\times2$ 分块形式的表示是
\begin{equation}
  \gamma^0 = \begin{pmatrix}0&1\\1&0\end{pmatrix},\qquad
  \gamma^i = \begin{pmatrix}0&\sigma^i\\ -\sigma^i&0\end{pmatrix}.
  \label{eq:3.26}
\end{equation}
这个表示称为 \emph{外尔或手征表示}。我们将发现它是一个特别方便的选择,并且全书将只用这个表示。(不过要小心,因为许多场论教科书选择不同的表示,其中 $\gamma^0$ 是对角的。此外,使用手征表示的书籍经常做出不同的符号约定选择。)

在我们的表示中,洛伦兹加速和旋转生成元为
\begin{equation}
  S^{0i} = \frac{i}{2}\begin{pmatrix}-\sigma^i&0\\0&\sigma^i\end{pmatrix}
  \label{eq:3.27}
\end{equation}
和
\begin{equation}
  S^{ij} = \frac{1}{2}\epsilon^{ijk}
    \begin{pmatrix}\sigma^k&0\\0&\sigma^k\end{pmatrix}.
  \label{eq:3.28}
\end{equation}
按 \eqref{eq:3.27} 和 \eqref{eq:3.28} 在洛伦兹加速和旋转下变换的四分量场 $\psi$ 称为 \emph{狄拉克旋量}。注意旋转生成元 $S^{ij}$ 正是三维旋量变换矩阵 \eqref{eq:3.25} 复制两次。洛伦兹加速生成元 $S^{0i}$ 不是厄米的,因此我们实现的洛伦兹加速不是幺正的(矢量表示 \eqref{eq:3.18} 也是如此)。事实上,洛伦兹群因为是``非紧致''的,没有忠实的、幺正的有限维表示。但这与我们无关,因为 $\psi$ 不是波函数;它是一个经典场。

现在我们有了 $\psi$ 的变换律,应该寻找一个合适的场方程。一种可能性就是克莱因-戈登方程:
\begin{equation}
  (\partial^2 + m^2)\psi = 0 .
  \label{eq:3.29}
\end{equation}
这之所以可行,是因为旋量变换矩阵 \eqref{eq:3.27} 和 \eqref{eq:3.28} 只作用于``内''空间;它们直接穿过微分算符。但我们可以写出一个更强的一阶方程,它蕴含 \eqref{eq:3.29} 但包含更多的信息。要做到这一点,我们需要知道 $\gamma$ 矩阵的另一个性质。经过简短的计算,你可以验证
\begin{equation}
  \{\gamma^\mu, S^{\nu\sigma}\} = (\mathcal{J}^{\nu\sigma})^{\mu}_{\ \rho}
    \gamma^\rho,
  \label{eq:3.30}
\end{equation}
或者等价地,
\begin{equation}
  \Bigl( 1 + \tfrac{i}{2}\omega_{\rho\sigma}S^{\rho\sigma} \Bigr)
    \gamma^\mu
    \Bigl( 1 - \tfrac{i}{2}\omega_{\rho\sigma}S^{\rho\sigma} \Bigr)
  = \Bigl( 1 - \tfrac{i}{2}\omega_{\rho\sigma}\mathcal{J}^{\rho\sigma}
    \Bigr)^{\mu}_{\ \nu}\gamma^\nu .
\end{equation}
这个方程正是
\begin{equation}
  \Lambda^{-1}_{1/2}\gamma^\mu\Lambda_{1/2}
    = \Lambda^{\mu}_{\ \nu}\gamma^\nu,
  \label{eq:3.31}
\end{equation}
的无穷小形式,其中
\begin{equation}
  \Lambda_{1/2} = \exp\Bigl( -\frac{i}{2}\Omega_{\mu\nu}S^{\mu\nu} \Bigr)
  \label{eq:3.32}
\end{equation}
是洛伦兹变换 $\Lambda$ 的旋量表示(比较 \eqref{eq:3.19})。式 \eqref{eq:3.31} 说明,$\gamma$ 矩阵在它们的矢量指标和旋量指标同时旋转下是不变的(就像 $\sigma^i$ 在空间旋转下不变一样)。换句话说,我们可以``认真对待 $\gamma^\mu$ 上的矢量指标 $\mu$'',把 $\gamma^\mu$ 与 $\partial_\mu$ 点乘,构成一个洛伦兹不变的微分算符。

我们现在准备好写下狄拉克方程。它就在这里:
\begin{equation}
  (i\gamma^\mu\partial_\mu - m)\psi(x) = 0 .
  \label{eq:3.33}
\end{equation}
为了证明它是洛伦兹不变的,写下左边的洛伦兹变换版本并计算:
\begin{align}
  (i\gamma^\mu\partial_\mu - m)\psi(x)
  &\to \bigl[ i\gamma^\mu(\Lambda^{-1})^{\nu}_{\ \mu}\partial_\nu - m \bigr]
    \Lambda_{1/2}\psi(\Lambda^{-1}x)\\
  &= \Lambda_{1/2}\bigl[ i\Lambda_{1/2}^{-1}\gamma^\mu\Lambda_{1/2}
    (\Lambda^{-1})^{\nu}_{\ \mu}\partial_\nu - m \bigr]\psi(\Lambda^{-1}x)\\
  &= \Lambda_{1/2}\bigl[ i\Lambda^{\mu}_{\ \rho}
    (\Lambda^{-1})^{\nu}_{\ \mu}\gamma^\rho\partial_\nu - m \bigr]
    \psi(\Lambda^{-1}x)\\
  &= \Lambda_{1/2}\bigl[ i\gamma^\rho\partial_\rho - m \bigr]
    \psi(\Lambda^{-1}x) = 0 .
\end{align}
为了看到狄拉克方程蕴含克莱因-戈登方程,用 $(-i\gamma^\nu\partial_\nu - m)$ 作用左边:
\begin{align}
  0 &= (-i\gamma^\nu\partial_\nu - m)(i\gamma^\mu\partial_\mu - m)\psi\\
  &= \Bigl( \gamma^\mu\gamma^\nu\partial_\mu\partial_\nu + m^2 \Bigr)\psi
    = \Bigl( \tfrac{1}{2}\{\gamma^\mu,\gamma^\nu\}\partial_\mu\partial_\nu
      + m^2 \Bigr)\psi
  = (\partial^2 + m^2)\psi .
\end{align}

为了写出狄拉克理论的拉格朗日量,我们必须弄清楚如何把两个狄拉克旋量相乘构成一个洛伦兹标量。明显的猜测 $\psi^\dagger\psi$ 并不奏效。在洛伦兹加速下,它变成 $\psi^\dagger\Lambda_{1/2}^\dagger\Lambda_{1/2}\psi$;如果洛伦兹加速矩阵是幺正的,我们就会有 $\Lambda_{1/2}^\dagger = \Lambda_{1/2}^{-1}$,一切都没问题。但 $\Lambda_{1/2}$ 不是幺正的,因为生成元 \eqref{eq:3.27} 不是厄米的。解决办法是定义
\begin{equation}
  \bar{\psi} = \psi^\dagger\gamma^0 .
  \label{eq:3.34}
\end{equation}
在由 $\Omega_{\mu\nu}$ 参数化的无穷小洛伦兹变换下,我们有
$\bar{\psi} \to \psi^\dagger\bigl( 1 + \frac{i}{2}\Omega_{\mu\nu}
(S^{\mu\nu})^\dagger \bigr)\gamma^0$。对 $\mu$ 和 $\nu$ 的求和有六个不同的非零项。在旋转项中,$\mu$ 和 $\nu$ 都非零,$(S^{\mu\nu})^\dagger = S^{\mu\nu}$ 且与 $\gamma^0$ 对易。在洛伦兹加速项中,$\mu$ 或 $\nu$ 为 0,$(S^{0i})^\dagger = -S^{0i}$ 但与 $\gamma^0$ 反对易。把 $\gamma^0$ 移到左边因此去掉了 $S^{\mu\nu}$ 上的剑号,给出变换律
\begin{equation}
  \bar{\psi} \to \bar{\psi}\Lambda_{1/2}^{-1},
  \label{eq:3.35}
\end{equation}
因此量 $\bar{\psi}\psi$ 是洛伦兹标量。类似地,你可以(借助 \eqref{eq:3.31})证明 $\bar{\psi}\gamma^\mu\psi$ 是洛伦兹矢量。

因此,正确的洛伦兹不变狄拉克拉格朗日量为
\begin{equation}
  \mathcal{L}_{\mathrm{Dirac}} = \bar{\psi}(i\gamma^\mu\partial_\mu
    - m)\psi .
  \label{eq:3.36}
\end{equation}
对 $\bar{\psi}$(或 $\psi^\dagger$)的欧拉-拉格朗日方程立即给出 \eqref{eq:3.33} 形式的狄拉克方程;对 $\psi$ 的欧拉-拉格朗日方程给出相同的方程,但取厄米共轭形式:
\begin{equation}
  i\partial_\mu\bar{\psi}\gamma^\mu + m\bar{\psi} = 0 .
  \label{eq:3.37}
\end{equation}

\subsection{外尔旋量}

从生成元 \eqref{eq:3.27} 和 \eqref{eq:3.28} 的分块对角形式可以明显看出,洛伦兹群的狄拉克表示是可约的。通过分别考虑每个分块,并写成
\begin{equation}
  \psi = \begin{pmatrix}\psi_L\\ \psi_R\end{pmatrix}.
  \label{eq:3.38}
\end{equation}
我们可以形成两个二维表示。二分量对象 $\psi_L$ 和 $\psi_R$ 称为 \emph{左手}和\emph{右手外尔旋量}。你可以容易地验证,在无穷小旋转 $\bm{\theta}$ 和洛伦兹加速 $\bm{\beta}$ 下,它们的变换律为
\begin{align}
  \psi_L &\to \Bigl( 1 - i\bm{\theta}\cdot\tfrac{\bm{\sigma}}{2}
    - \bm{\beta}\cdot\tfrac{\bm{\sigma}}{2} \Bigr)\psi_L,\\
  \psi_R &\to \Bigl( 1 - i\bm{\theta}\cdot\tfrac{\bm{\sigma}}{2}
    + \bm{\beta}\cdot\tfrac{\bm{\sigma}}{2} \Bigr)\psi_R .
  \label{eq:3.39}
\end{align}
这些变换律通过复共轭相联系;利用恒等式
\begin{equation}
  \sigma^2\bm{\sigma}^{*} = -\bm{\sigma}\sigma^2,
  \label{eq:3.40}
\end{equation}
不难证明量 $\sigma^2\psi_L^{*}$ 像右手旋量那样变换。

用 $\psi_L$ 和 $\psi_R$ 表示,狄拉克方程为
\begin{equation}
  \begin{pmatrix}
    -m & i\sigma^\mu\partial_\mu\\
    i\bar{\sigma}^\mu\partial_\mu & -m
  \end{pmatrix}
  \begin{pmatrix}\psi_L\\ \psi_R\end{pmatrix} = 0 .
  \label{eq:3.41}
\end{equation}
两个洛伦兹群表示 $\psi_L$ 和 $\psi_R$ 被狄拉克方程中的质量项混合。但如果我们令 $m=0$,$\psi_L$ 和 $\psi_R$ 的方程就解耦:
\begin{align}
  i(\partial_0 - \bm{\sigma}\cdot\nabla)\psi_L &= 0;\\
  i(\partial_0 + \bm{\sigma}\cdot\nabla)\psi_R &= 0 .
  \label{eq:3.42}
\end{align}
这些方程称为 \emph{外尔方程};它们在处理中微子和弱相互作用理论时特别重要。

可以把记号稍微整理一下。定义
\begin{equation}
  \sigma^\mu = (1, \bm{\sigma}),\qquad
  \bar{\sigma}^\mu = (1, -\bm{\sigma}),
  \label{eq:3.43}
\end{equation}
使得
\begin{equation}
  \gamma^\mu = \begin{pmatrix}
    0 & \sigma^\mu\\
    \bar{\sigma}^\mu & 0
  \end{pmatrix}.
  \label{eq:3.44}
\end{equation}
($\sigma$ 上的横杠与 $\psi$ 上的横杠完全无关。)于是狄拉克方程可以写为
\begin{equation}
  \begin{pmatrix}
    -m & i\sigma\cdot\partial\\
    i\bar{\sigma}\cdot\partial & -m
  \end{pmatrix}
  \begin{pmatrix}\psi_L\\ \psi_R\end{pmatrix} = 0
  \label{eq:3.45}
\end{equation}
外尔方程变为
\begin{equation}
  i\bar{\sigma}\cdot\partial\,\psi_L = 0;\qquad
  i\sigma\cdot\partial\,\psi_R = 0 .
  \label{eq:3.46}
\end{equation}

\section{狄拉克方程的自由粒子解}

为了对狄拉克方程的物理有一些感觉,现在让我们讨论它的平面波解。由于狄拉克场 $\psi$ 满足克莱因-戈登方程,我们立即知道它可以写成平面波的线性组合:
\begin{equation}
  \psi(x) = u(p)e^{-ip\cdot x},\qquad \text{其中 } p^2 = m^2 .
  \label{eq:3.47}
\end{equation}
眼下我们将集中关注正频解,即 $p^0>0$。列矢量 $u(p)$ 必须满足一个额外的约束,这是把 \eqref{eq:3.47} 代入狄拉克方程得到的:
\begin{equation}
  (i\gamma^\mu\partial_\mu - m)u(p)e^{-ip\cdot x} = 0
  \qquad\Rightarrow\qquad
  (\gamma^\mu p_\mu - m)u(p) = 0 .
  \label{eq:3.48}
\end{equation}
在静止系中分析这个方程最容易,其中 $p = p_0 = (m, \mathbf{0})$;一般 $p$ 的解可以通过用 $\Lambda_{1/2}$ 洛伦兹加速得到。在静止系中,式 \eqref{eq:3.48} 变为
\begin{equation}
  (m\gamma^0 - m)u(p_0) = m
    \begin{pmatrix}-1&1\\1&-1\end{pmatrix}u(p_0) = 0 ,
\end{equation}
解为
\begin{equation}
  u(p_0) = \sqrt{m}\begin{pmatrix}\xi\\ \xi\end{pmatrix},
  \label{eq:3.49}
\end{equation}
对任意数值二分量旋量 $\xi$。我们按惯例把 $\xi$ 归一化为 $\xi^\dagger\xi = 1$;因子 $\sqrt{m}$ 是为了将来方便而插入的。我们可以通过看旋转生成元 \eqref{eq:3.28} 来解释旋量 $\xi$:$\xi$ 在旋转下作为旋转群的一个普通二分量旋量变换,因此以通常的方式决定狄拉克解的自旋取向。例如,当 $\xi = \binom{1}{0}$ 时,粒子具有沿 3 方向向上的自旋。

注意,在应用狄拉克方程之后,我们只能自由选择 $u(p)$ 四个分量中的两个。这正是我们想要的,因为自旋 $1/2$ 粒子只有两个物理态——自旋向上和自旋向下。(当然,我们谈论粒子和自旋有点为时过早。当我们将在第 3.5 节量子化狄拉克理论时,会证明狄拉克粒子的自旋角动量为 $\hbar/2$;眼下,只需注意到对任何动量 $p$ 都有两个可能的解 $u(p)$。)

既然我们在静止系中有了 $u(p)$ 的一般形式,就可以通过洛伦兹加速得到任何其他参考系中的 $u(p)$。考虑沿 3 方向的洛伦兹加速。首先我们应该提醒自己洛伦兹加速对 4-动量矢量做了什么。在无穷小形式中,
\begin{equation}
  p^\mu \to (1 - i\eta K^3)^{\mu}_{\ \nu}p^\nu
    = \begin{pmatrix}
        \cosh\eta & \sinh\eta\\
        \sinh\eta & \cosh\eta
      \end{pmatrix}
      \begin{pmatrix}m\\ 0\end{pmatrix}
    = \begin{pmatrix}m\cosh\eta\\ m\sinh\eta\end{pmatrix},
\end{equation}
其中 $\eta$ 是某个无穷小参数。对有限的 $\eta$,我们必须写
\begin{equation}
  p^\mu = \begin{pmatrix}
    \cosh\eta & \sinh\eta\\
    \sinh\eta & \cosh\eta
  \end{pmatrix}
  \begin{pmatrix}m\\ 0\end{pmatrix}
  = \begin{pmatrix}
    E\\ p^3
  \end{pmatrix},
  \label{eq:3.50}
\end{equation}
其中 $E = m\cosh\eta$, $p^3 = m\sinh\eta$。参数 $\eta$ 称为 \emph{快度}。它是在相继洛伦兹加速下可相加的量。

现在对 $u(p)$ 应用同样的洛伦兹加速。根据式 \eqref{eq:3.27} 和 \eqref{eq:3.32},
\begin{align}
  u(p) &= \exp\Bigl( \tfrac{1}{2}\eta\gamma^0\gamma^3 \Bigr)
    \sqrt{m}\begin{pmatrix}\xi\\ \xi\end{pmatrix}
  = \begin{pmatrix}
      e^{\eta/2} & 0\\
      0 & e^{-\eta/2}
    \end{pmatrix}
    \sqrt{m}\begin{pmatrix}\xi\\ \xi\end{pmatrix}\\
  &= \sqrt{m}
    \begin{pmatrix}
      e^{\eta/2}\xi\\
      e^{-\eta/2}\xi
    \end{pmatrix}
  = \begin{pmatrix}
      \sqrt{p\cdot\sigma}\,\xi\\
      \sqrt{p\cdot\bar{\sigma}}\,\xi
    \end{pmatrix}.
  \label{eq:3.51}
\end{align}
最后一行可以化简给出
\begin{equation}
  u(p) = \begin{pmatrix}
    \sqrt{p\cdot\sigma}\,\xi\\
    \sqrt{p\cdot\bar{\sigma}}\,\xi
  \end{pmatrix},
  \label{eq:3.52}
\end{equation}
其中约定,在取矩阵的平方根时,我们对每个本征值取正根。$u(p)$ 的这个表达式不仅更紧凑,而且对任意方向的 $\mathbf{p}$ 也成立。在处理这种形式的表达式时,知道下面的恒等式常常很有用
\begin{equation}
  \sqrt{p\cdot\sigma}\,\sqrt{p\cdot\bar{\sigma}} = m
  \qquad\text{(把 $\sigma$ 和 $\bar{\sigma}$ 互换也同样成立)。}
  \label{eq:3.53}
\end{equation}

对 $\xi = \binom{1}{0}$(自旋沿 3 轴向上),我们有
\begin{equation}
  u(p) = \begin{pmatrix}
    \sqrt{p^0-p^3}\,\binom{1}{0}\\
    \sqrt{p^0+p^3}\,\binom{1}{0}
  \end{pmatrix}
  = \begin{pmatrix}
    \sqrt{E-p^3}\binom{1}{0}\\
    \sqrt{E+p^3}\binom{1}{0}
  \end{pmatrix},
  \label{eq:3.54}
\end{equation}
而对 $\xi = \binom{0}{1}$(自旋沿 3 轴向下),我们有
\begin{equation}
  u(p) = \begin{pmatrix}
    \sqrt{p\cdot\sigma}\,\binom{0}{1}\\
    \sqrt{p\cdot\bar{\sigma}}\,\binom{0}{1}
  \end{pmatrix}
  = \begin{pmatrix}
    \sqrt{E+p^3}\binom{0}{1}\\
    \sqrt{E-p^3}\binom{0}{1}
  \end{pmatrix}.
  \label{eq:3.55}
\end{equation}
在极限 $\eta\to\infty$ 下,这些态退化为无质量粒子的二分量旋量。(我们现在看到了 \eqref{eq:3.49} 中因子 $\sqrt{m}$ 的原因:它使旋量表达式在无质量极限下保持有限。)

解 \eqref{eq:3.54} 和 \eqref{eq:3.55} 是螺旋度算符
\begin{equation}
  h = \frac{1}{2}\hat{\mathbf{p}}\cdot\bm{\sigma}
  = \frac{1}{2}\hat{\mathbf{p}}\cdot\mathbf{J}.
  \label{eq:3.56}
\end{equation}
的本征态。$h = +1/2$ 的粒子称为右手粒子,$h = -1/2$ 的粒子称为左手粒子。有质量粒子的螺旋度依赖于参考系,因为人们总可以洛伦兹加速到一个其动量方向相反的参考系(但它的自旋不变)。对以光速运动的无质量粒子,人们无法进行这样的洛伦兹加速。

无质量粒子在螺旋度本征态下 $u(p)$ 的极其简单的形式,使得这种粒子的行为易于理解。在第 1 章中,它使我们能够在无质量极限下猜测 $e^{+}e^{-}\to\mu^{+}\mu^{-}$ 散射截面的形式。在随后的章节中,我们经常会先做一次不动脑筋的计算,然后在高能极限下看螺旋度本征态,以理解我们所做的一切。

顺便提一下,我们现在准备好理解外尔旋量记号 $\psi_L$ 和 $\psi_R$ 的由来。外尔方程的解是具有确定螺旋度的态,分别对应左手和右手粒子。螺旋度的洛伦兹不变性(对无质量粒子)在外尔旋量的记号中显而易见,因为 $\psi_L$ 和 $\psi_R$ 生活在洛伦兹群的不同表示中。

方便起见,把 $u(p)$ 的归一化条件写成洛伦兹不变的形式。上面我们看到 $u^\dagger u$ 不是洛伦兹不变的。类似地,
\begin{equation}
  u^\dagger(p)u(p) = 2E_p\,\xi^\dagger\xi .
  \label{eq:3.57}
\end{equation}
为了构造洛伦兹标量,我们定义
\begin{equation}
  \bar{u}(p) = u^\dagger(p)\gamma^0 .
  \label{eq:3.58}
\end{equation}
于是通过几乎相同的计算,
\begin{equation}
  \bar{u}u = 2m\,\xi^\dagger\xi .
  \label{eq:3.59}
\end{equation}
一旦我们还要求二分量旋量 $\xi$ 按通常方式归一化:$\xi^\dagger\xi = 1$,这就成为我们的归一化条件。也常选择正交的基旋量 $\xi^1$ 和 $\xi^2$(如 $\binom{1}{0}$ 和 $\binom{0}{1}$)。对无质量粒子,式 \eqref{eq:3.59} 是平凡的,所以我们必须把归一化条件写成 \eqref{eq:3.57} 的形式。

让我们总结到目前为止的讨论。狄拉克方程的一般解可以写成平面波的线性组合。正频波具有形式
\begin{equation}
  \psi(x) = u(p)e^{-ip\cdot x},\qquad p^2 = m^2,\quad p^0>0 .
  \label{eq:3.60}
\end{equation}
$u(p)$ 有两个线性无关的解,
\begin{equation}
  u^s(p),\qquad s = 1, 2,
  \label{eq:3.61}
\end{equation}
我们按
\begin{equation}
  \bar{u}^r(p)u^s(p) = 2m\delta^{rs}
  \qquad\text{或}\qquad
  u^{r\dagger}(p)u^s(p) = 2E_p\delta^{rs}.
  \label{eq:3.62}
\end{equation}
归一化。以完全相同的方式,我们可以找到负频解:
\begin{equation}
  \psi(x) = v(p)e^{+ip\cdot x},\qquad p^2 = m^2,\quad p^0>0 .
  \label{eq:3.63}
\end{equation}
(注意,我们选择把 $+$ 号放进指数中,而不是让 $p^0<0$。) $v(p)$ 有两个线性无关的解,
\begin{equation}
  v^s(p),\qquad s = 1, 2,
  \label{eq:3.64}
\end{equation}
其中 $\eta^s$ 是二分量旋量的另一组基。这些解按
\begin{equation}
  \bar{v}^r(p)v^s(p) = -2m\delta^{rs}
  \qquad\text{或}\qquad
  v^{r\dagger}(p)v^s(p) = +2E_p\delta^{rs}.
  \label{eq:3.65}
\end{equation}
归一化。$u$ 和 $v$ 也彼此正交:
\begin{equation}
  \bar{u}^r(p)v^s(p) = \bar{v}^r(p)u^s(p) = 0 .
  \label{eq:3.66}
\end{equation}
要小心,因为 $u^{r\dagger}(p)v^s(p)\neq 0$ 且 $v^{r\dagger}(p)u^s(p)\neq 0$。不过,注意
\begin{equation}
  u^{r\dagger}(p)v^s(-\mathbf{p}) =
  v^{r\dagger}(-\mathbf{p})u^s(p) = 0 ,
  \label{eq:3.67}
\end{equation}
其中我们在每个旋量乘积的一个因子中改变了 3-动量的符号。

\subsection{自旋求和}

在计算费曼图时,我们经常希望对费米子的极化态求和。通过一个简单的计算,我们可以推导出相关的完备性关系:
\begin{align}
  \sum_{s=1,2} u^s(p)\bar{u}^s(p) &=
    \begin{pmatrix}
      \sqrt{p\cdot\sigma}\,\xi^s\\
      \sqrt{p\cdot\bar{\sigma}}\,\xi^s
    \end{pmatrix}
    \Bigl( (\xi^{s\dagger}\sqrt{p\cdot\bar{\sigma}})
      \quad (\xi^{s\dagger}\sqrt{p\cdot\sigma}) \Bigr)\\
  &= \begin{pmatrix}
      \sqrt{p\cdot\sigma}\,(\sum_s \xi^s\xi^{s\dagger})\sqrt{p\cdot\bar{\sigma}}
        & \sqrt{p\cdot\sigma}\,(\sum_s \xi^s\xi^{s\dagger})\sqrt{p\cdot\sigma}\\
      \sqrt{p\cdot\bar{\sigma}}\,(\sum_s \xi^s\xi^{s\dagger})
        \sqrt{p\cdot\bar{\sigma}}
        & \sqrt{p\cdot\bar{\sigma}}\,(\sum_s \xi^s\xi^{s\dagger})
          \sqrt{p\cdot\sigma}
    \end{pmatrix}\\
  &= \begin{pmatrix}
      \sqrt{p\cdot\sigma}\,\mathbf{1}\sqrt{p\cdot\bar{\sigma}}
        & \sqrt{p\cdot\sigma}\,\mathbf{1}\sqrt{p\cdot\sigma}\\
      \sqrt{p\cdot\bar{\sigma}}\,\mathbf{1}\sqrt{p\cdot\bar{\sigma}}
        & \sqrt{p\cdot\bar{\sigma}}\,\mathbf{1}\sqrt{p\cdot\sigma}
    \end{pmatrix}
  = \begin{pmatrix}
      m & p\cdot\sigma\\
      p\cdot\bar{\sigma} & m
    \end{pmatrix}.
  \label{eq:3.68}
\end{align}
在第二行我们用了
\begin{equation}
  \sum_{s=1,2}\xi^s\xi^{s\dagger} = \begin{pmatrix}1&0\\0&1\end{pmatrix}.
\end{equation}
于是我们得到想要的公式,
\begin{equation}
  \sum_{s} u^s(p)\bar{u}^s(p) = \gamma\cdot p + m .
  \label{eq:3.69}
\end{equation}
类似地,
\begin{equation}
  \sum_{s} v^s(p)\bar{v}^s(p) = \gamma\cdot p - m .
  \label{eq:3.70}
\end{equation}
组合 $\gamma\cdot p$ 出现得如此频繁,以至于 Feynman 引入了记号
\begin{equation}
  \slashed{p} = \gamma^\mu p_\mu .
  \label{eq:3.71}
\end{equation}
从现在起我们将经常使用这个记号。

\section{狄拉克矩阵与狄拉克场双线性量}

我们在第 3.2 节看到量 $\bar{\psi}\psi$ 是洛伦兹标量。也容易证明 $\bar{\psi}\gamma^\mu\psi$ 是 4-矢量——我们在写下狄拉克拉格朗日量 \eqref{eq:3.36} 时用到了这一事实。现在让我们问一个更一般的问题:考虑表达式 $\bar{\psi}\Gamma\psi$,其中 $\Gamma$ 是任意 $4\times4$ 常数矩阵。我们能否把这个表达式分解成在洛伦兹群下具有确定变换性质的项?答案是可以,只要我们把 $\Gamma$ 用下面的十六个 $4\times4$ 矩阵基写出,这些矩阵定义为 $\gamma$ 矩阵的反对称组合:
\begin{equation}
  \begin{array}{lll}
    1 & \text{(1 个)} & \text{标量}\\
    \gamma^\mu & \text{(4 个)} & \text{矢量}\\
    \sigma^{\mu\nu} = \tfrac{i}{2}[\gamma^\mu,\gamma^\nu] &
      \text{(6 个)} & \text{张量}\\
    \gamma^\mu\gamma^5 & \text{(4 个)} & \text{赝矢量}\\
    \gamma^5 & \text{(1 个)} & \text{赝标量}
  \end{array}
\end{equation}
这些矩阵的洛伦兹变换性质容易确定。例如,
\begin{align}
  \bar{\psi}[\gamma^\mu,\gamma^\nu]\psi
  &\to (\psi^\dagger\Lambda_{1/2}^\dagger)
    \tfrac{i}{2}[\gamma^\mu,\gamma^\nu](\Lambda_{1/2}\psi)\\
  &= \tfrac{i}{2}\Lambda^{\mu}_{\ \alpha}\Lambda^{\nu}_{\ \beta}
    \psi^\dagger[\gamma^\alpha,\gamma^\beta]\psi .
\end{align}
每一组矩阵都作为逐次升高秩的反对称张量变换。

最后两组矩阵可以通过引入一个额外的伽马矩阵来化简,
\begin{equation}
  \gamma^5 = i\gamma^0\gamma^1\gamma^2\gamma^3
    = -\frac{i}{4!}\epsilon^{\mu\nu\rho\sigma}
      \gamma_\mu\gamma_\nu\gamma_\rho\gamma_\sigma .
  \label{eq:3.72}
\end{equation}
于是 $\sigma^{\mu\nu}\gamma^5 = -i\epsilon^{\mu\nu\rho\sigma}\sigma_{\rho\sigma}$,且 $\gamma^\mu\gamma^5 = +i\epsilon^{\mu\nu\rho\sigma}\sigma_{\nu\rho}\gamma_\sigma$。矩阵 $\gamma^5$ 具有下列性质,所有这些性质都可以用 \eqref{eq:3.72} 和反对易关系 \eqref{eq:3.22b} 验证:
\begin{align}
  (\gamma^5)^\dagger &= \gamma^5; \label{eq:3.73}\\
  (\gamma^5)^2 &= 1; \label{eq:3.74}\\
  \{\gamma^5, \gamma^\mu\} &= 0 . \label{eq:3.75}
\end{align}
最后这个性质意味着 $[\gamma^5, S^{\mu\nu}] = 0$。因此狄拉克表示必定是可约的,因为本征值不同的 $\gamma^5$ 本征矢量变换时互不混合(这个可约性判据称为舒尔引理)。在我们的基中,
\begin{equation}
  \gamma^5 = \begin{pmatrix}-1&0\\0&1\end{pmatrix}
  \label{eq:3.76}
\end{equation}
以分块对角形式。所以只有左手(右手)分量的狄拉克旋量是 $\gamma^5$ 的、本征值为 $-1$($+1$)的本征态,而且这些旋量确实变换时互不混合,正如我们在第 3.2 节显式看到的那样。

现在让我们重写 $4\times4$ 矩阵的表,并引入一些标准术语:
\begin{equation}
  \begin{array}{lll}
    \text{标量} & 1 & \text{(1)}\\
    \text{矢量} & \gamma^\mu & \text{(4)}\\
    \text{张量} & \sigma^{\mu\nu} = \tfrac{i}{2}[\gamma^\mu,\gamma^\nu]
      & \text{(6)}\\
    \text{赝矢量} & \gamma^\mu\gamma^5 & \text{(4)}\\
    \text{赝标量} & \gamma^5 & \text{(1)}
  \end{array}
  \label{eq:3.77}
\end{equation}
赝矢量和赝标量这两个术语源于如下事实:这些量在连续洛伦兹变换下分别作为矢量和标量变换,但在宇称变换下会多一个符号变化(我们将在第 3.6 节讨论)。

从矢量和赝矢量矩阵出发,我们可以由狄拉克场双线性量构成两个流:
\begin{equation}
  j^\mu(x) = \bar{\psi}(x)\gamma^\mu\psi(x);\qquad
  j^{\mu 5}(x) = \bar{\psi}(x)\gamma^\mu\gamma^5\psi(x) .
  \label{eq:3.78}
\end{equation}
让我们计算这些流的散度,假设 $\psi$ 满足狄拉克方程:
\begin{align}
  \partial_\mu j^\mu &= (\partial_\mu\bar{\psi})\gamma^\mu\psi
    + \bar{\psi}\gamma^\mu(\partial_\mu\psi)\\
  &= (im\bar{\psi})\psi + \bar{\psi}(-im\psi) = 0 .
  \label{eq:3.79}
\end{align}
因此,如果 $\psi(x)$ 满足狄拉克方程,$j^\mu$ 总是守恒的。当我们把狄拉克场与电磁场耦合时,$j^\mu$ 将变成电流密度。类似地,可以计算
\begin{equation}
  \partial_\mu j^{\mu 5} = 2im\bar{\psi}\gamma^5\psi .
  \label{eq:3.80}
\end{equation}
如果 $m=0$,这个流(常称为轴矢流)也是守恒的。于是构造线性组合是有用的
\begin{equation}
  j_L^\mu = \bar{\psi}\gamma^\mu\frac{1-\gamma^5}{2}\psi
    = \bar{\psi}_L\gamma^\mu\psi_L;\qquad
  j_R^\mu = \bar{\psi}\gamma^\mu\frac{1+\gamma^5}{2}\psi
    = \bar{\psi}_R\gamma^\mu\psi_R .
  \label{eq:3.81}
\end{equation}
当 $m=0$ 时,它们分别是左手粒子和右手粒子的电流密度,并且分别守恒。

两个流 $j^\mu(x)$ 和 $j^{\mu 5}(x)$ 是对应于两个变换
\begin{equation}
  \psi(x) \to e^{i\alpha}\psi(x)
  \qquad\text{以及}\qquad
  \psi(x) \to e^{i\alpha\gamma^5}\psi(x) .
\end{equation}
的诺特流。其中第一个是狄拉克拉格朗日量 \eqref{eq:3.36} 的对称性。第二个,称为 \emph{手征变换},是 $\mathcal{L}$ 中导数项的对称性,但不是质量项的对称性;因此,诺特定理确认轴矢流只有在 $m=0$ 时才守恒。

狄拉克双线性量的乘积服从交换关系,称为费尔兹恒等式。我们只讨论其中最简单的一个,它在本书后面会多次用到。这个最简单的恒等式用式 \eqref{eq:3.38} 引入的二分量外尔旋量来写最为容易。该关系的核心是式 \eqref{eq:3.43} 定义的 $2\times2$ 矩阵 $\sigma^\mu$ 的恒等式:
\begin{equation}
  (\sigma^\mu)_{\alpha\dot{\alpha}}(\sigma_\mu)_{\beta\dot{\beta}}
    = 2\epsilon_{\alpha\beta}\epsilon_{\dot{\alpha}\dot{\beta}} .
  \label{eq:3.82}
\end{equation}
(这里 $\alpha, \beta$ 等是旋量指标,$\epsilon$ 是反对称符号。)可以通过如下观察来理解这个关系:指标 $\alpha, \gamma$ 在 $\psi_L$ 的洛伦兹表示中变换,而 $\beta, \delta$ 在 $\psi_R$ 的独立表示中变换,整个量必须是洛伦兹不变量。或者,也可以直接显式验证 \eqref{eq:3.82} 的 16 个分量。

通过把恒等式 \eqref{eq:3.82} 夹在狄拉克旋量 $u_1, u_2, u_3, u_4$ 的右手部分(即下半部分)之间,我们发现恒等式
\begin{equation}
  (\bar{u}_{1R}\sigma^\mu u_{2R})(\bar{u}_{3R}\sigma_\mu u_{4R})
    = 2\epsilon_{\alpha\beta}\,
    (\bar{u}_{1R})_\alpha(\bar{u}_{3R})_\beta\,
    \epsilon_{\gamma\delta}\,
    (u_{2R})_\gamma(u_{4R})_\delta ,
  \label{eq:3.83}
\end{equation}
它是标准的费尔兹恒等式之一。

\section{狄拉克场的量子化}

我们现在准备量子化狄拉克场,遵循我们在第 2.3 节对克莱因-戈登场使用的同一程序。我们将把经典场 $\psi$ 及其共轭动量提升为服从正则对易关系的算符。然而,对狄拉克场,我们必须施加\emph{反对易}关系,以便理论只有正能量激发。这个程序的形式论证——自旋-统计定理——将在下面讨论,遵循上一节引导我们走向费米统计的同一逻辑。

狄拉克理论的经典拉格朗日量为
\begin{equation}
  \mathcal{L} = \bar{\psi}(i\gamma^\mu\partial_\mu - m)\psi .
  \label{eq:3.84}
\end{equation}
$\psi_\alpha$ 的共轭动量为
\begin{equation}
  \pi_\alpha = \frac{\partial\mathcal{L}}{\partial\dot{\psi}_\alpha}
    = i(\psi^\dagger\gamma^0\gamma^0)_\alpha = i\psi^\dagger_\alpha .
  \label{eq:3.85}
\end{equation}
哈密顿密度为
\begin{equation}
  \mathcal{H} = \pi\dot{\psi} - \mathcal{L}
    = \bar{\psi}(-i\gamma^i\partial_i + m)\psi .
  \label{eq:3.86}
\end{equation}

我们试图通过施加正则\emph{对易}关系来量子化这个理论,
\begin{equation}
  [\psi_\alpha(\mathbf{x}), \psi_\beta^\dagger(\mathbf{y})]
    = \delta^{(3)}(\mathbf{x}-\mathbf{y})\delta_{\alpha\beta};
  \qquad
  [\psi_\alpha(\mathbf{x}), \psi_\beta(\mathbf{y})]
  = [\psi_\alpha^\dagger(\mathbf{x}), \psi_\beta^\dagger(\mathbf{y})] = 0 .
  \label{eq:3.87}
\end{equation}
我们把 $\psi(x)$ 和 $\bar{\psi}(x)$ 用产生和湮灭算符 $a$, $a^\dagger$, $b$, $b^\dagger$ 展开为傅里叶积分:
\begin{equation}
  \psi(x) = \int \frac{d^3p}{(2\pi)^3}\ \frac{1}{\sqrt{2E_{\mathbf{p}}}}
    \sum_s \Bigl( a_{\mathbf{p}}^s u^s(p)e^{-ip\cdot x}
      + b_{\mathbf{p}}^{s\dagger} v^s(p)e^{+ip\cdot x} \Bigr);
  \label{eq:3.88}
\end{equation}
\begin{equation}
  \bar{\psi}(x) = \int \frac{d^3p}{(2\pi)^3}\ \frac{1}{\sqrt{2E_{\mathbf{p}}}}
    \sum_s \Bigl( b_{\mathbf{p}}^s \bar{v}^s(p)e^{-ip\cdot x}
      + a_{\mathbf{p}}^{s\dagger} \bar{u}^s(p)e^{+ip\cdot x} \Bigr).
  \label{eq:3.89}
\end{equation}
正如对克莱因-戈登场一样,我们发现对易关系 \eqref{eq:3.87} 意味着
\begin{equation}
  \{a_{\mathbf{p}}^r, a_{\mathbf{q}}^{s\dagger}\}
  = \{b_{\mathbf{p}}^r, b_{\mathbf{q}}^{s\dagger}\}
    = (2\pi)^3\delta^{(3)}(\mathbf{p}-\mathbf{q})\delta^{rs},
  \label{eq:3.90}
\end{equation}
而所有其他对易子为零,前提是场满足狄拉克方程。这是一个会带来麻烦的对易关系,正如我们现在将要展示的那样。

让我们考虑传播振幅 $\langle 0|\bar{\psi}(y)\psi(x)|0\rangle$ 以及与之对应的 $\langle 0|\psi(x)\bar{\psi}(y)|0\rangle$。第一个振幅只从 $\psi(x)$ 中带 $b^\dagger$ 的项和 $\bar{\psi}(y)$ 中带 $b$ 的项获得贡献;第二个只从带 $a$ 和 $a^\dagger$ 的项获得贡献。但对类空间隔,因果性要求这两个振幅抵消:
\begin{equation}
  \langle 0|\psi_a(x)\bar{\psi}_b(y)|0\rangle
  = \langle 0|\bar{\psi}_b(y)\psi_a(x)|0\rangle
  \qquad\text{对 } (x-y)^2<0 .
  \label{eq:3.91}
\end{equation}
让我们计算这些振幅,使用假定的对易关系 \eqref{eq:3.87}。例如,我们发现
\begin{equation}
  \langle 0|\psi_a(x)\bar{\psi}_b(y)|0\rangle =
    \int \frac{d^3p}{(2\pi)^3}\ \frac{1}{2E_{\mathbf{p}}}
    \sum_s u_a^s(p)\bar{u}_b^s(p)\, e^{-ip\cdot(x-y)},
  \label{eq:3.92}
\end{equation}
和
\begin{equation}
  \langle 0|\bar{\psi}_b(y)\psi_a(x)|0\rangle =
    \int \frac{d^3p}{(2\pi)^3}\ \frac{1}{2E_{\mathbf{p}}}
    \sum_s v_a^s(p)\bar{v}_b^s(p)\, e^{+ip\cdot(x-y)},
  \label{eq:3.93}
\end{equation}
正如对克莱因-戈登场一样。但在克莱因-戈登情形中,我们分别从这两部分各得到对易子的一项:粒子从 $y$ 到 $x$ 的传播被反粒子从 $x$ 到 $y$、在光锥外的传播抵消。这里两项都来自第一部分 $\langle 0|\bar{\psi}(x)\psi(y)|0\rangle$,因为第二部分为零。抵消发生在正能量粒子与负能量粒子之间,两者都从 $y$ 传播到 $x$。

这个观察实际上可以引导我们找到负能量问题的解决之道。我们在量子化狄拉克理论时所做的假设中,必有一个是不正确的。因此让我们忘掉假定的对易关系 \eqref{eq:3.87} 和 \eqref{eq:3.90},看看我们能否找到一种方式,使正能量粒子能向两个方向传播。我们也将不得不放弃把真空 $|0\rangle$ 定义为被所有 $a_{\mathbf{p}}^s$ 和 $b_{\mathbf{p}}^s$ 湮灭的态。不过,我们将保留 $\psi(x)$ 和 $\bar{\psi}(x)$ 作为海森堡算符的表达式 \eqref{eq:3.88} 和 \eqref{eq:3.89},因为如果 $\psi(x)$ 和 $\bar{\psi}(x)$ 求解狄拉克方程,它们必定可以分解成这样的平面波解。

首先考虑传播振幅 $\langle 0|\psi(x)\bar{\psi}(y)|0\rangle$,它要表示一个从 $y$ 传播到 $x$ 的正能量粒子。在这种情况下,我们希望(海森堡)态 $\bar{\psi}(y)|0\rangle$ 只由正能量或负频分量组成(因为海森堡态 $\psi_H = e^{iHt}\psi_S$)。因此 $\bar{\psi}(y)$ 中只有 $a^\dagger$ 项可以贡献,这意味着 $a_{\mathbf{p}}^s$ 必须湮灭真空。类似地,$\langle 0|\psi(x)$ 只能包含正频分量。于是我们有
\begin{equation}
  \langle 0|\psi_a(x)\bar{\psi}_b(y)|0\rangle
    = \int \frac{d^3p}{(2\pi)^3}\ \frac{1}{2E_{\mathbf{p}}}
    \sum_s u_a^s(p)\bar{u}_b^s(p)\, e^{-ip\cdot(x-y)}
    \langle 0|a_{\mathbf{p}}^s a_{\mathbf{p}}^{s\dagger}|0\rangle .
  \label{eq:3.94}
\end{equation}
即使不知道如何交换 $a_{\mathbf{p}}^r$ 和 $a_{\mathbf{q}}^{s\dagger}$,我们也可以利用平移和旋转不变性,对矩阵元 $\langle 0|a_{\mathbf{p}}^r a_{\mathbf{q}}^{s\dagger}|0\rangle$ 说些什么。如果基态 $|0\rangle$ 要在平移下不变,我们必须有
$|0\rangle = e^{iP\cdot x}|0\rangle$。此外,由于 $a_{\mathbf{p}}^{r\dagger}$ 产生动量 $q$,我们可以用式 \eqref{eq:2.48} 计算
\begin{equation}
  \langle 0|a_{\mathbf{p}}^r a_{\mathbf{q}}^{s\dagger}|0\rangle
    = \langle 0|a_{\mathbf{p}}^r a_{\mathbf{q}}^{s\dagger}
      e^{iP\cdot x}|0\rangle
    = e^{i(p-q)\cdot x}\langle 0|a_{\mathbf{p}}^r a_{\mathbf{q}}^{s\dagger}|0\rangle .
\end{equation}
这说明,如果 $\langle 0|a_{\mathbf{p}}^r a_{\mathbf{q}}^{s\dagger}|0\rangle$ 非零,则必须有 $p = q$。类似地,可以证明 $|0\rangle$ 的旋转不变性意味着 $r = s$。(这应当是直观清楚的,并且可以在本节后面讨论角动量算符之后加以检验。)从这些考虑我们得出结论,矩阵元可以写成
\begin{equation}
  \langle 0|a_{\mathbf{p}}^r a_{\mathbf{q}}^{s\dagger}|0\rangle
    = (2\pi)^3\delta^{(3)}(\mathbf{p}-\mathbf{q})\delta^{rs} \cdot A(p),
\end{equation}
其中 $A(p)$ 迄今未定。不过注意,如果态的范数总是正的(正如在任何合格的希尔伯特空间中都应当如此),$A(p)$ 必须大于零。我们现在可以回到 \eqref{eq:3.94},写出
\begin{equation}
  \langle 0|\psi_a(x)\bar{\psi}_b(y)|0\rangle
    = (i\slashed{\partial}_x + m)_{ab}\,
      \int \frac{d^3p}{(2\pi)^3}\ \frac{1}{2E_{\mathbf{p}}}\,
      e^{-ip\cdot(x-y)}\, A(p) .
  \label{eq:3.95}
\end{equation}
这个表达式只有在 $A(p)$ 是洛伦兹标量,即 $A(p) = A(p^2)$ 时,才在洛伦兹加速下正确地保持不变。由于 $p^2 = m^2$,$A$ 必须是一个常数。所以最终我们得到
\begin{equation}
  \langle 0|\psi_a(x)\bar{\psi}_b(y)|0\rangle
    = (i\slashed{\partial}_x + m)_{ab}\,
      \int \frac{d^3p}{(2\pi)^3}\ \frac{1}{2E_{\mathbf{p}}}\,
      e^{-ip\cdot(x-y)} \cdot A .
  \label{eq:3.96}
\end{equation}
类似地,在振幅 $\langle 0|\bar{\psi}(y)\psi(x)|0\rangle$ 中,我们希望唯一贡献来自 $\bar{\psi}(y)$ 的正频项和 $\psi(x)$ 的负频项。所以 $a^\dagger$ 仍然湮灭真空,但 $b^\dagger$ 不湮灭真空。于是通过与上面给出的论证完全相同的论证,我们有
\begin{equation}
  \langle 0|\bar{\psi}_b(y)\psi_a(x)|0\rangle
    = -(i\slashed{\partial}_x + m)_{ab}\,
      \int \frac{d^3p}{(2\pi)^3}\ \frac{1}{2E_{\mathbf{p}}}\,
      e^{+ip\cdot(x-y)} \cdot B ,
  \label{eq:3.97}
\end{equation}
其中 $B$ 是另一个正常数。负号很重要;它来自 $v^s$ 的完备性关系 \eqref{eq:3.70} 和指数因子中 $x$ 的符号。它意味着我们不可能在光锥外有 $\langle 0|[\psi(x),\bar{\psi}(y)]|0\rangle = 0$:如果 $A = -B$,两项 \eqref{eq:3.96} 和 \eqref{eq:3.97} 确实会抵消,但这是不可能的,因为 $A$ 和 $B$ 都必须为正。

不过,解决办法现在已经近在眼前。通过令 $A = B = 1$,容易得到(在光锥外)
\begin{equation}
  \langle 0|\psi_a(x)\bar{\psi}_b(y)|0\rangle
    = -\langle 0|\bar{\psi}_b(y)\psi_a(x)|0\rangle .
\end{equation}
也就是说,旋量场在类空间隔处\emph{反对易}。这足以保持因果性,因为所有合理的可观测量(如能量、电荷和粒子数)都由偶数个旋量场构成;对任何这样的可观测量 $\mathcal{O}_1$ 和 $\mathcal{O}_2$,对 $(x-y)^2<0$,我们仍然有 $[\mathcal{O}_1(x), \mathcal{O}_2(y)] = 0$。而且值得注意的是,为狄拉克场设定反对易关系解决了负能量问题。等时反对易关系将为
\begin{equation}
  \{\psi_a(\mathbf{x}), \psi_b^\dagger(\mathbf{y})\}
    = \delta^{(3)}(\mathbf{x}-\mathbf{y})\delta_{ab};
  \qquad
  \{\psi_a(\mathbf{x}), \psi_b(\mathbf{y})\}
  = \{\psi_a^\dagger(\mathbf{x}), \psi_b^\dagger(\mathbf{y})\} = 0 .
  \label{eq:3.98}
\end{equation}
我们可以像以前一样(式 \eqref{eq:3.88})用 $a_{\mathbf{p}}^s$ 和 $b_{\mathbf{p}}^{s\dagger}$ 展开 $\psi(x)$。产生和湮灭算符现在必须服从
\begin{equation}
  \{a_{\mathbf{p}}^r, a_{\mathbf{q}}^{s\dagger}\}
  = \{b_{\mathbf{p}}^r, b_{\mathbf{q}}^{s\dagger}\}
    = (2\pi)^3\delta^{(3)}(\mathbf{p}-\mathbf{q})\delta^{rs}
  \label{eq:3.99}
\end{equation}
(所有其他反对易子为零),以便 \eqref{eq:3.98} 被满足。

另一次计算给出哈密顿量,
\begin{equation}
  H = \int \frac{d^3p}{(2\pi)^3}\ E_{\mathbf{p}}
    \sum_s \bigl( a_{\mathbf{p}}^{s\dagger}a_{\mathbf{p}}^s
      - b_{\mathbf{p}}^s b_{\mathbf{p}}^{s\dagger} \bigr),
\end{equation}
它和以前一样;$b_{\mathbf{p}}^{s\dagger}$ 仍然产生负能量。然而,关系
$\{b_{\mathbf{p}}^r, b_{\mathbf{q}}^{s\dagger}\}
= (2\pi)^3\delta^{(3)}(\mathbf{p}-\mathbf{q})\delta^{rs}$ 在 $b_{\mathbf{p}}^r$ 与 $b_{\mathbf{q}}^{s\dagger}$ 之间是对称的。所以让我们简单地重新定义
\begin{equation}
  b_{\mathbf{p}}^s = B_{\mathbf{p}}^s,\qquad
  b_{\mathbf{p}}^{s\dagger} = B_{\mathbf{p}}^{s\dagger} .
  \label{eq:3.100}
\end{equation}
它们当然服从完全相同的反对易关系,但现在哈密顿量中的第二项为
\begin{equation}
  -E_{\mathbf{p}}b_{\mathbf{p}}^s b_{\mathbf{p}}^{s\dagger}
    = +E_{\mathbf{p}}b_{\mathbf{p}}^{s\dagger}b_{\mathbf{p}}^s
      - (\text{const}).
\end{equation}
如果我们选择 $|0\rangle$ 为被 $a_{\mathbf{p}}^s$ 和 $b_{\mathbf{p}}^s$ 湮灭的态,那么 $|0\rangle$ 的所有激发都具有正能量。

发生了什么?为了更好地理解这个技巧,让我们暂时放下场论,考虑一个只有一对 $b$ 和 $b^\dagger$ 算符的理论,它们服从 $\{b,b^\dagger\} = 1$ 和 $\{b,b\} = \{b^\dagger,b^\dagger\} = 0$。选择一个满足 $b|0\rangle = 0$ 的态 $|0\rangle$。那么 $b^\dagger|0\rangle$ 是一个新态;称它为 $|1\rangle$。这个态满足 $b|1\rangle = |0\rangle$ 和 $b^\dagger|1\rangle = 0$。所以 $b$ 和 $b^\dagger$ 作用在一个只有两个态 $|0\rangle$ 和 $|1\rangle$ 的希尔伯特空间上。我们或许会说,$|0\rangle$ 表示一个``空''态,$b^\dagger$ ``填充''这个态。但我们同样可以把 $|1\rangle$ 称为空态,并说 $b$ —— 填充它。这两种描述完全等价,直到我们指定某个可观测量,使我们在物理上能够区分这些态。在我们的情形中,正确的选择是取能量较低的态为空态。而把剑号放在产生正能量的算符上不那么令人困惑。这正是我们所做的。

顺便注意,由于 $(b^\dagger)^2 = 0$,这个态不能被填充两次。更一般地,反对易关系意味着任何多粒子态在两个粒子交换下是反对称的:$a_{\mathbf{p}}^{s\dagger}a_{\mathbf{q}}^{r\dagger}|0\rangle =
-a_{\mathbf{q}}^{r\dagger}a_{\mathbf{p}}^{s\dagger}|0\rangle$。因此我们得出结论:如果阶梯算符服从反对易关系,相应的粒子就服从费米-狄拉克统计。

我们刚刚证明,为了确保真空只有正能量激发,我们必须用反对易关系来量子化狄拉克场;在这些条件下,与狄拉克场相联系的粒子服从费米-狄拉克统计。这个结论是一个更一般结果的一部分,该结果首先由 Pauli 导出。%
\footnote{W.~Pauli, \emph{Phys.\ Rev.} \textbf{58}, 716 (1940), 转载于 Schwinger (1958)。严格的论述见 R.~F.~Streater 和 A.~S.~Wightman, \emph{PCT, Spin and Statistics, and All That} (Benjamin/Cummings, Reading, Mass., 1964)。}:洛伦兹不变性、正能量、正范数与因果性共同意味着整数自旋的粒子服从玻色-爱因斯坦统计,而半奇整数自旋的粒子服从费米-狄拉克统计。

\subsection{量子化的狄拉克场}

让我们以系统的方式总结量子化狄拉克理论的结果。既然尘埃已经落定,我们应该整理一下记号:从现在起,我们把 $b_{\mathbf{p}}^s$(降低态能量的算符)就写为 $b_{\mathbf{p}}^s$,把 $b_{\mathbf{p}}^{s\dagger}$ 写为 $b_{\mathbf{p}}^{s\dagger}$。我们以后需要用的所有表达式都列在下面;上面那些与之不同的对应表达式应当忘掉。

首先我们写出场算符:
\begin{equation}
  \psi(x) = \int \frac{d^3p}{(2\pi)^3}\ \frac{1}{\sqrt{2E_{\mathbf{p}}}}
    \sum_s \Bigl( a_{\mathbf{p}}^s u^s(p)e^{-ip\cdot x}
      + b_{\mathbf{p}}^{s\dagger} v^s(p)e^{+ip\cdot x} \Bigr);
  \label{eq:3.101}
\end{equation}
\begin{equation}
  \bar{\psi}(x) = \int \frac{d^3p}{(2\pi)^3}\ \frac{1}{\sqrt{2E_{\mathbf{p}}}}
    \sum_s \Bigl( b_{\mathbf{p}}^s \bar{v}^s(p)e^{-ip\cdot x}
      + a_{\mathbf{p}}^{s\dagger} \bar{u}^s(p)e^{+ip\cdot x} \Bigr).
  \label{eq:3.102}
\end{equation}
产生和湮灭算符服从反对易规则
\begin{equation}
  \{a_{\mathbf{p}}^r, a_{\mathbf{q}}^{s\dagger}\}
  = \{b_{\mathbf{p}}^r, b_{\mathbf{q}}^{s\dagger}\}
    = (2\pi)^3\delta^{(3)}(\mathbf{p}-\mathbf{q})\delta^{rs},
  \label{eq:3.103}
\end{equation}
所有其他反对易子为零。$\psi$ 和 $\psi^\dagger$ 的等时反对易关系于是为
\begin{equation}
  \{\psi_a(\mathbf{x}), \psi_b^\dagger(\mathbf{y})\}
    = \delta^{(3)}(\mathbf{x}-\mathbf{y})\delta_{ab};
  \qquad
  \{\psi_a(\mathbf{x}), \psi_b(\mathbf{y})\}
  = \{\psi_a^\dagger(\mathbf{x}), \psi_b^\dagger(\mathbf{y})\} = 0 .
  \label{eq:3.104}
\end{equation}
真空 $|0\rangle$ 被定义为满足
\begin{equation}
  a_{\mathbf{p}}^s|0\rangle = b_{\mathbf{p}}^s|0\rangle = 0 .
  \label{eq:3.105}
\end{equation}
的态。哈密顿量可以写成
\begin{equation}
  H = \int \frac{d^3p}{(2\pi)^3}\ \sum_s E_{\mathbf{p}}
    \Bigl( a_{\mathbf{p}}^{s\dagger}a_{\mathbf{p}}^s
      + b_{\mathbf{p}}^{s\dagger}b_{\mathbf{p}}^s \Bigr),
  \label{eq:3.106}
\end{equation}
其中我们略去了由 $b_{\mathbf{p}}^s$ 和 $b_{\mathbf{p}}^{s\dagger}$ 反对易带来的无穷大常数项。由此我们看到,真空就是能量最低的态,这正是我们所希望的。动量算符为
\begin{equation}
  \mathbf{P} = \int d^3x\ \psi^\dagger(-i\nabla)\psi
    = \int \frac{d^3p}{(2\pi)^3}\ \sum_s \mathbf{p}
      \bigl( a_{\mathbf{p}}^{s\dagger}a_{\mathbf{p}}^s
        + b_{\mathbf{p}}^{s\dagger}b_{\mathbf{p}}^s \bigr).
  \label{eq:3.107}
\end{equation}
因此,$a_{\mathbf{p}}^{s\dagger}$ 和 $b_{\mathbf{p}}^{s\dagger}$ 都产生具有能量 $+E_{\mathbf{p}}$ 和动量 $\mathbf{p}$ 的粒子。我们将把 $a_{\mathbf{p}}^{s\dagger}$ 产生的粒子称为 \emph{费米子},把 $b_{\mathbf{p}}^{s\dagger}$ 产生的粒子称为 \emph{反费米子}。

单粒子态
\begin{equation}
  |\mathbf{p}, s\rangle = \sqrt{2E_{\mathbf{p}}}\,a_{\mathbf{p}}^{s\dagger}|0\rangle
  \label{eq:3.108}
\end{equation}
的定义使得它们的内积
\begin{equation}
  \langle\mathbf{p}, r|\mathbf{q}, s\rangle
    = 2E_{\mathbf{p}}\,(2\pi)^3\delta^{(3)}(\mathbf{p}-\mathbf{q})\delta^{rs}
  \label{eq:3.109}
\end{equation}
是洛伦兹不变的。这意味着在希尔伯特空间的态上实现洛伦兹变换的算符 $U(\Lambda)$ 是幺正的,尽管对洛伦兹加速而言 $\Lambda_{1/2}$ 不是幺正的。

做一次自洽性检验,看看 $U(\Lambda)$ 是否在 $\psi(x)$ 上实现了正确的变换,这将令人安心。因此计算
\begin{equation}
  U\psi(x)U^{-1} = U \int \frac{d^3p}{(2\pi)^3}\ \frac{1}{\sqrt{2E_{\mathbf{p}}}}
    \sum_s \Bigl( a_{\mathbf{p}}^s u^s(p)e^{-ip\cdot x}
      + b_{\mathbf{p}}^{s\dagger} v^s(p)e^{ip\cdot x} \Bigr) U^{-1}.
  \label{eq:3.110}
\end{equation}
我们可以集中于第一项;第二项完全类似。式 \eqref{eq:3.108} 意味着 $a_{\mathbf{p}}^{s\dagger}$ 按
\begin{equation}
  U(\Lambda)a_{\mathbf{p}}^{s\dagger}U^{-1}(\Lambda)
    = a_{\Lambda\mathbf{p}}^{s\dagger},
  \label{eq:3.111}
\end{equation}
变换,假设自旋量子化轴平行于洛伦兹加速轴或旋转轴。为了利用这个关系来计算 \eqref{eq:3.110},把积分改写为
\begin{equation}
  \int \frac{d^3p}{(2\pi)^3}\ \frac{1}{\sqrt{2E_{\mathbf{p}}}}\ a_{\mathbf{p}}^s
    = \int \frac{d^3p}{(2\pi)^3}\ \frac{1}{\sqrt{2E_{\mathbf{p}}}}\,
      u^s(p)\, e^{-ip\cdot x}
    = \int \frac{d^3p}{(2\pi)^3}\ \frac{1}{2E_{\mathbf{p}}}\,
      u^s(p)\, e^{-ip\cdot x}\, \sqrt{2E_{\mathbf{p}}} .
\end{equation}
第二个因子被 $U$ 以简单的方式变换,而第一个是洛伦兹不变的积分。因此,如果我们应用 \eqref{eq:3.111} 并作代换 $p = \Lambda p$,式 \eqref{eq:3.110} 变为
\begin{equation}
  U(\Lambda)\psi(x)U^{-1}(\Lambda) =
    \int \frac{d^3p}{(2\pi)^3}\ \frac{1}{\sqrt{2E_{\mathbf{p}}}}\,
    \sum_s a_{\mathbf{p}}^{s\dagger}
    u^s(\Lambda^{-1}p)\, e^{-ip\cdot \Lambda x} + \cdots .
\end{equation}
但 $u^s(\Lambda^{-1}p) = \Lambda_{1/2}u^s(p)$,所以我们确实有
\begin{equation}
  U(\Lambda)\psi(x)U^{-1}(\Lambda) = \Lambda_{1/2}\psi(\Lambda x),
  \label{eq:3.112}
\end{equation}
正如所宣称的那样。

在狄拉克理论中还有一个重要的守恒量。在第 3.4 节我们看到了流 $j^\mu = \bar{\psi}\gamma^\mu\psi$ 是守恒的。与此流相联系的电荷是
\begin{equation}
  Q = \int d^3x\ \psi^\dagger(x)\psi(x)
    = \int \frac{d^3p}{(2\pi)^3}\ \sum_s
      \bigl( a_{\mathbf{p}}^{s\dagger}a_{\mathbf{p}}^s
        - b_{\mathbf{p}}^{s\dagger}b_{\mathbf{p}}^s \bigr),
  \label{eq:3.113}
\end{equation}
或者,如果我们忽略另一个无穷大常数,
\begin{equation}
  Q = \int \frac{d^3p}{(2\pi)^3}\ \sum_s
    \bigl( a_{\mathbf{p}}^{s\dagger}a_{\mathbf{p}}^s
      - b_{\mathbf{p}}^s b_{\mathbf{p}}^{s\dagger} \bigr)
    = \int \frac{d^3p}{(2\pi)^3}\ \sum_s
      \bigl( a_{\mathbf{p}}^{s\dagger}a_{\mathbf{p}}^s
        - b_{\mathbf{p}}^{s\dagger}b_{\mathbf{p}}^s \bigr).
  \label{eq:3.114}
\end{equation}
因此 $a_{\mathbf{p}}^{s\dagger}$ 产生电荷为 $+1$ 的费米子,而 $b_{\mathbf{p}}^{s\dagger}$ 产生电荷为 $-1$ 的反费米子。当我们把狄拉克场与电磁场耦合时,将会看到 $Q$ 正是电荷(相差一个依赖于我们想要描述哪种粒子的常数因子;例如,对电子而言,电荷为 $Qe$)。

在量子电动力学中,我们将用旋量场 $\psi$ 来描述电子和正电子。$a_{\mathbf{p}}^{s\dagger}$ 产生的粒子是电子;它们具有能量 $E_{\mathbf{p}}$、动量 $\mathbf{p}$、自旋 $1/2$(极化与 $u^s$ 相应),电荷 $+1$(以 $e$ 为单位)。$b_{\mathbf{p}}^{s\dagger}$ 产生的粒子是正电子;它们具有能量 $E_{\mathbf{p}}$、动量 $\mathbf{p}$、自旋 $1/2$(极化与 $u^s$ 相反),电荷 $-1$。态 $\psi_a(x)|0\rangle$ 在位置 $x$ 含有一个正电子,其极化对应于所选的旋量分量。类似地,$\bar{\psi}_a(x)|0\rangle$ 是在位置 $x$ 含有一个电子的态。

\subsection{狄拉克传播子}

计算狄拉克场的传播振幅现在已经是直截了当的练习:
\begin{align}
  \langle 0|\psi_a(x)\bar{\psi}_b(y)|0\rangle
  &= \int \frac{d^3p}{(2\pi)^3}\ \frac{1}{2E_{\mathbf{p}}}
    (i\slashed{\partial}_x + m)_{ab}\, e^{-ip\cdot(x-y)}
  = (i\slashed{\partial}_x + m)_{ab}\,
    \int \frac{d^3p}{(2\pi)^3}\ \frac{1}{2E_{\mathbf{p}}}\,
    e^{-ip\cdot(x-y)} ,
  \label{eq:3.115}\\
  \langle 0|\bar{\psi}_b(y)\psi_a(x)|0\rangle
  &= (i\slashed{\partial}_x + m)_{ab}\,
    \int \frac{d^3p}{(2\pi)^3}\ \frac{1}{2E_{\mathbf{p}}}\,
    e^{+ip\cdot(x-y)} .
  \label{eq:3.116}
\end{align}
正如我们对克莱因-戈登方程所做的那样,我们可以为狄拉克方程构造满足各种边界条件的格林函数。例如,推迟格林函数为
\begin{equation}
  S_R(x-y) = \theta(x^0-y^0)\ \langle 0|\{\psi_a(x), \bar{\psi}_b(y)\}|0\rangle .
  \label{eq:3.117}
\end{equation}
容易验证
\begin{equation}
  S_R(x-y) = (i\slashed{\partial}_x + m)\,D_R(x-y),
  \label{eq:3.118}
\end{equation}
因为在右边涉及 $\partial_0\theta(x^0-y^0)$ 的项为零。利用 \eqref{eq:3.118} 以及 $\slashed{\partial}\slashed{\partial} = \partial^2$ 这一事实,我们看到 $S_R$ 是狄拉克算符的格林函数:
\begin{equation}
  (i\slashed{\partial}_x - m)\,S_R(x-y) = i\delta^{(4)}(x-y)\cdot
    \mathbb{1}_{4\times4} .
  \label{eq:3.119}
\end{equation}
狄拉克算符的格林函数也可以通过傅里叶变换求得。把 $S_R(x-y)$ 展开为傅里叶积分,并用 $(i\slashed{\partial}_x - m)$ 作用两边,我们得到
\begin{equation}
  iS_R(x-y) = \int \frac{d^4p}{(2\pi)^4}\ (\slashed{p} - m)\,
    e^{-ip\cdot(x-y)}\, S_R(p),
  \label{eq:3.120}
\end{equation}
因此
\begin{equation}
  S_R(p) = \frac{i}{\slashed{p} - m}
    = \frac{i(\slashed{p} + m)}{p^2 - m^2} .
  \label{eq:3.121}
\end{equation}
为了得到推迟格林函数,我们必须沿第 30 页所示的围道计算 \eqref{eq:3.121} 中的 $p^0$ 积分。对 $x^0>y^0$,我们在下方闭合围道,取到两个极点,得到 \eqref{eq:3.115} 和 \eqref{eq:3.116} 之和。对 $x^0<y^0$,我们在上方闭合围道,得到零。

具有费曼边界条件的格林函数由第 31 页所示的围道定义:
\begin{equation}
  S_F(x-y) = \int \frac{d^4p}{(2\pi)^4}\
    \frac{i(\slashed{p} + m)}{p^2 - m^2 + i\epsilon}\ e^{-ip\cdot(x-y)}
  = \begin{cases}
      \langle 0|\psi(x)\bar{\psi}(y)|0\rangle & \text{对 } x^0>y^0
        \text{ (在下方闭合围道)}\\
      -\langle 0|\bar{\psi}(y)\psi(x)|0\rangle & \text{对 } x^0<y^0
        \text{ (在上方闭合围道)}
    \end{cases}
  = \langle 0|T\psi(x)\bar{\psi}(y)|0\rangle ,
  \label{eq:3.122}
\end{equation}
其中我们选择把旋量场的编时乘积定义为在算符交换时带有额外负号。这个负号在费米子的量子场论中极其重要;我们将在第 4.7 节再次遇到它。

与克莱因-戈登理论一样,费曼传播子的表达式 \eqref{eq:3.122} 是本章最有用的结果。当我们用费曼图做微扰计算时,将把因子 $S_F(p)$ 与每条内部费米子线联系起来。

\section{狄拉克理论的离散对称性}

在上一节中,我们讨论了在狄拉克理论的希尔伯特空间上实现连续洛伦兹变换。我们发现,对每个变换 $\Lambda$ 都有一个幺正算符 $U(\Lambda)$,它在场上诱导出正确的变换:
\begin{equation}
  U(\Lambda)\psi(x)U^{-1}(\Lambda) = \Lambda_{1/2}\psi(\Lambda x).
  \label{eq:3.123}
\end{equation}
在本节中,我们将讨论在狄拉克场上实现各种分立对称性的类似算符。

除了连续洛伦兹变换之外,还有两种其他的时空操作是拉格朗日量潜在的对称性:宇称和时间反演。宇称记为 $P$,把 $(t,\mathbf{x})$ 送到 $(t,-\mathbf{x})$,反转空间的手性。时间反演记为 $T$,把 $(t,\mathbf{x})$ 送到 $(-t,\mathbf{x})$,交换前后光锥。这两种操作都不能通过从恒等元出发的连续洛伦兹变换来实现。然而,两者都保持闵可夫斯基间隔 $x^2 = t^2 - \mathbf{x}^2$。按照标准术语,连续洛伦兹变换被称为固有正时洛伦兹群 $L_+$。于是,完整的洛伦兹群分裂成四个不相连的子集,如下所示:
\begin{equation}
  \begin{array}{ll}
    L_+^\uparrow = L_+ & \text{``固有''}\\
    L_- = PL_+ & \text{``非固有''}\\
    L_+^\downarrow = TL_+ & \text{``非正时''}\\
    L_-^\uparrow = PL_+^\downarrow = TL_+^\uparrow &
      \text{``非固有非正时''}
  \end{array}
\end{equation}
在我们讨论 $P$ 和 $T$ 的同时,方便起见还要讨论第三种(非时空的)分立操作:电荷共轭,记为 $C$。在此操作下,粒子和反粒子互换。

虽然任何相对论性场论在 $L_+^\uparrow$ 下都必须不变,但它不必在 $P$, $T$ 或 $C$ 下不变。这些对称操作在现实世界中处于什么地位?从实验我们知道,自然界三种力——引力、电磁和强相互作用——对 $P$, $C$ 和 $T$ 是对称的。弱相互作用分别破坏 $C$ 和 $P$,但保持 $CP$ 和 $T$。但某些稀有过程(迄今观察到的都涉及中性 $K$ 介子)也显示出 $CP$ 和 $T$ 破坏。所有观测都表明,组合 $CPT$ 是自然界的一个完美对称性。

当前被接受的弱相互作用理论模型是格拉肖-温伯格-萨拉姆规范理论,在第 20 章中描述。这个理论以最强的方式破坏 $C$ 和 $P$。$C$ 和 $P$ 在最容易观察到的过程中碰巧是相当好的对称性,这实际上是一个意外(虽然不是完全的偶然)。另一方面,没有人知道一个真正优美的破坏 $CP$ 的理论。在当前理论中,当存在三代(或更多)费米子时,有一个参数的余地,如果该参数非零,就会引起 $CP$ 破坏。但这个参数的值并不比电子质量的值更容易理解;$CP$ 破坏的物理起源仍然是一个谜。我们将在第 20.3 节进一步讨论这个问题。

\subsection{宇称}

有了这个引言,现在让我们讨论 $P$, $T$ 和 $C$ 对狄拉克粒子和场的作用。首先考虑宇称。算符 $P$ 应当反转粒子的动量而不翻转其自旋:
\begin{equation}
  \mathbf{p} \to -\mathbf{p}\qquad \text{(镜像)。}
\end{equation}
数学上,这意味着 $P$ 应当由一个幺正算符(严格说应称为 $U(P)$,但我们只称它为 $P$)实现,例如,它把态 $a_{\mathbf{p}}^{s\dagger}|0\rangle$ 变换为 $a_{-\mathbf{p}}^{s\dagger}|0\rangle$。换句话说,我们希望
\begin{equation}
  P a_{\mathbf{p}}^s P = \eta_a a_{-\mathbf{p}}^s
  \qquad\text{以及}\qquad
  P b_{\mathbf{p}}^s P = \eta_b b_{-\mathbf{p}}^s,
  \label{eq:3.124}
\end{equation}
其中 $\eta_a$ 和 $\eta_b$ 是可能的相位。这些相位受到如下条件的限制:宇称算符作用两次应当使可观测量回到其原始值。由于可观测量由偶数个费米子算符构成,这就要求 $\eta_a^2, \eta_b^2 = \pm 1$。

正如连续洛伦兹变换在狄拉克场上由 $4\times4$ 常数矩阵 $\Lambda_{1/2}$ 实现一样,宇称变换也应当由一个 $4\times4$ 常数矩阵表示。为了找到这个矩阵并确定 $\eta_a$ 和 $\eta_b$,我们计算 $P$ 对 $\psi(x)$ 的作用。利用 \eqref{eq:3.124},我们有
\begin{equation}
  P\psi(x)P = \int \frac{d^3p}{(2\pi)^3}\ \frac{1}{\sqrt{2E_{\mathbf{p}}}}
    \sum_s \Bigl( \eta_a a_{-\mathbf{p}}^s u^s(p)e^{-ip\cdot x}
      + \eta_b^{*} b_{\mathbf{p}}^{s\dagger} v^s(p)e^{ip\cdot x} \Bigr).
  \label{eq:3.125}
\end{equation}
现在把变量变换为 $\mathbf{p}' = (p^0, -\mathbf{p})$。注意
$p\cdot x = p'\cdot(t,-\mathbf{x})$。还有
\begin{equation}
  u^s(p) = \gamma^0 u^s(p'),\qquad
  v^s(p) = -\gamma^0 v^s(p').
\end{equation}
所以 \eqref{eq:3.125} 变为
\begin{equation}
  P\psi(x)P = \int \frac{d^3p}{(2\pi)^3}\ \frac{1}{\sqrt{2E_{\mathbf{p}}}}
    \sum_s \Bigl( \eta_a a_{\mathbf{p}}^s \gamma^0 u^s(p)e^{-ip\cdot
      (t,-\mathbf{x})} - \eta_b^{*} b_{\mathbf{p}}^{s\dagger} \gamma^0 v^s(p)
      e^{ip\cdot(t,-\mathbf{x})} \Bigr).
\end{equation}
它应当等于某个常数矩阵乘以 $\psi(t,-\mathbf{x})$,而如果我们令 $\eta_b^{*} = -\eta_a$,确实如此。这意味着
\begin{equation}
  \eta_a\eta_b = -\eta_a^2 = -1 .
  \label{eq:3.126}
\end{equation}
于是我们得到了 $\psi(x)$ 宇称变换的最终形式,
\begin{equation}
  P\psi(t,\mathbf{x})P = \eta_a\gamma^0\psi(t,-\mathbf{x}) .
  \label{eq:3.127}
\end{equation}
了解各种狄拉克场双线性量在宇称下如何变换将非常重要(例如,在写下拉格朗日量时)。回忆五个双线性量是
\begin{equation}
  \bar{\psi}\psi,\qquad \bar{\psi}\gamma^\mu\psi,\qquad
  \bar{\psi}\sigma^{\mu\nu}\psi,\qquad
  \bar{\psi}\gamma^\mu\gamma^5\psi,\qquad
  \bar{\psi}i\gamma^5\psi .
\end{equation}
在宇称下,我们发现
\begin{align}
  P\bar{\psi}\psi\,P &= +\bar{\psi}\psi(t,-\mathbf{x});\\
  P\bar{\psi}\gamma^\mu\psi\,P &=
    +(-1)^\mu\bar{\psi}\gamma^\mu\psi(t,-\mathbf{x});\\
  P\bar{\psi}\sigma^{\mu\nu}\psi\,P &=
    (-1)^\mu(-1)^\nu\bar{\psi}\sigma^{\mu\nu}\psi(t,-\mathbf{x});\\
  P\bar{\psi}\gamma^\mu\gamma^5\psi\,P &=
    -(-1)^\mu\bar{\psi}\gamma^\mu\gamma^5\psi(t,-\mathbf{x});\\
  P\bar{\psi}i\gamma^5\psi\,P &=
    -\bar{\psi}i\gamma^5\psi(t,-\mathbf{x}),
\end{align}
其中对 $\mu = 0$, $(-1)^\mu = 1$;对 $\mu = 1, 2, 3$, $(-1)^\mu = -1$。

\subsection{时间反演}

现在我们转向更微妙的时间反演情形。我们不会描述原始文献中给出的完整处理;相反,我们将仅对结果作简要阐述,这有助于理解 $T$ 对狄拉克场的作用。

时间反演操作 $T$ 应当改变时间的方向,并且在某种意义上``反转''量子力学演化。要完整描述 $T$ 对量子力学系统的作用,我们需要把 $T$ 看作\emph{反线性}(或反幺正)算符。由于这可能是一个新概念,让我们简要解释一下。幺正算符 $U$ 作用于态的线性叠加为
\begin{equation}
  U(a|A\rangle + b|B\rangle) = aU|A\rangle + bU|B\rangle .
\end{equation}
另一方面,反线性算符 $\Theta$ 作用为
\begin{equation}
  \Theta(a|A\rangle + b|B\rangle) = a^{*}\Theta|A\rangle + b^{*}\Theta|B\rangle .
\end{equation}
时间反演是反线性的,因为例如可以预期,$\langle B|A\rangle$(一个复数)在 $T$ 下被变换为其复共轭:$T\langle B|A\rangle = \langle B|A\rangle^{*}$。此外,如果 $H$ 是一个时间反演不变系统的哈密顿量,使得 $T$ 与 $H$ 对易,那么时间演化算符 $e^{-iHt}$ 满足 $Te^{-iHt} = e^{iHt}T$,这正说明 $T$ 反转时间的方向。

方便起见,把狄拉克场的反线性变换 $T$ 参数化如下。由于 $T$ 反转动量和自旋,我们预期
\begin{equation}
  T\psi(t,\mathbf{x})T = \eta_T\, C\,\psi^{*}(-t,\mathbf{x})
\end{equation}
其中 $\eta_T$ 是某个相位,$C$ 是某个 $4\times4$ 矩阵(不要与电荷共轭算符混淆;这里 $C$ 被称为时间反演矩阵)。实际上,为了保持记号简单,让我们定义
\begin{equation}
  T\psi(t,\mathbf{x})T = \gamma^1\gamma^3\psi^{*}(-t,\mathbf{x})
  \label{eq:3.128}
\end{equation}
(一个整体相位已被略去)。利用式 \eqref{eq:3.67} 的记号,我们把费米子湮灭算符的时间反演变换定义如下:
\begin{equation}
  T a_{\mathbf{p}}^s T = a_{-\mathbf{p}}^s,\qquad
  T b_{\mathbf{p}}^s T = b_{-\mathbf{p}}^s .
  \label{eq:3.129}
\end{equation}
(额外的整体相位对我们其余的讨论没有影响,为简单起见略去。)关系 \eqref{eq:3.129} 使我们能够计算 $T$ 对费米子场 $\psi(x)$ 的作用:
\begin{align}
  T\psi(t,\mathbf{x})T &= \int \frac{d^3p}{(2\pi)^3}\ \frac{1}{\sqrt{2E_{\mathbf{p}}}}
    \sum_s \Bigl( a_{-\mathbf{p}}^s u^{s*}(p)e^{+ip\cdot x}
      + b_{-\mathbf{p}}^{s\dagger} v^{s*}(p)e^{-ip\cdot x} \Bigr)\\
  &= \gamma^1\gamma^3\,\psi(-t,\mathbf{x}) .
  \label{eq:3.130}
\end{align}
在最后一步我们用了 $p\cdot(t,-\mathbf{x}) = -p\cdot(-t,\mathbf{x})$。与宇称的情形一样,我们推导出了费米子场 $\psi(x)$ 的简单变换律。粒子和反粒子变换律中的相对负号在这里也存在,它隐含在 $v^{-s}$ 中两次翻转的旋量中。

现在我们可以检查 $T$ 对各种双线性量的作用。首先我们需要
\begin{equation}
  T\bar{\psi}T = (T\psi T)^{\dagger}\gamma^0
    = \psi^{\dagger}(-t,\mathbf{x})[\gamma^1\gamma^3]^{\dagger}\gamma^0
    = \bar{\psi}(-t,\mathbf{x})[-\gamma^1\gamma^3] .
  \label{eq:3.131}
\end{equation}
那么标量双线性量的变换为
\begin{equation}
  T\bar{\psi}\psi\,T = \bar{\psi}(-\gamma^1\gamma^3)(\gamma^1\gamma^3)\psi(-t,\mathbf{x})
    = \bar{\psi}\psi(-t,\mathbf{x}) .
  \label{eq:3.132}
\end{equation}
当 $T$ 穿过 $i$ 时,赝标量获得一个额外的负号:
\begin{equation}
  Ti\bar{\psi}\gamma^5\psi\,T = -i\bar{\psi}(-\gamma^1\gamma^3)\gamma^5
    (\gamma^1\gamma^3)\psi = -i\bar{\psi}\gamma^5\psi(-t,\mathbf{x}).
\end{equation}
对矢量,我们必须分别计算 $\mu = 0, 1, 2, 3$ 四种情形。经过一点计算,你应当得到
\begin{equation}
  T\bar{\psi}\gamma^\mu\psi\,T =
    \bar{\psi}(-\gamma^1\gamma^3)(\gamma^\mu)^{*}(\gamma^1\gamma^3)\psi
    = \begin{cases}
        +\bar{\psi}\gamma^\mu\psi(-t,\mathbf{x}) & \text{对 } \mu = 0;\\
        -\bar{\psi}\gamma^\mu\psi(-t,\mathbf{x}) & \text{对 } \mu = 1, 2, 3 .
      \end{cases}
  \label{eq:3.133}
\end{equation}
这正是我们想要的矢量(如流密度)变换性质。你可以验证赝矢量在时间反演下以完全相同的方式变换。

\subsection{电荷共轭}

三种分立对称性中最后一种是粒子-反粒子对称性 $C$。把 $C$ 实现为幺正线性算符不会有任何问题。电荷共轭按惯例定义为把具有给定自旋取向的费米子变成具有相同自旋取向的反费米子。因此,费米子湮灭算符的一个方便变换选择是
\begin{equation}
  C a_{\mathbf{p}}^s C = b_{\mathbf{p}}^s;\qquad
  C b_{\mathbf{p}}^s C = a_{\mathbf{p}}^s .
  \label{eq:3.134}
\end{equation}
同样,为简单起见我们忽略可能存在的额外相位。

接下来我们想要算出 $C$ 对 $\psi(x)$ 的作用。首先我们需要 $v^s(p)$ 和 $u^s(p)$ 之间的关系。利用 \eqref{eq:3.67} 和 \eqref{eq:3.55},
\begin{equation}
  v^s(p) = \begin{pmatrix}
    \sqrt{p\cdot\sigma}\,\eta^s\\
    \sqrt{p\cdot\bar{\sigma}}\,\eta^s
  \end{pmatrix}
  = \begin{pmatrix}
    -i\sigma^2\sqrt{p\cdot\bar{\sigma}}\,\xi^{s*}\\
    i\sigma^2\sqrt{p\cdot\sigma}\,\xi^{s*}
  \end{pmatrix}
  = -i\gamma^2\begin{pmatrix}
    \sqrt{p\cdot\sigma}\,\xi^{s*}\\
    \sqrt{p\cdot\bar{\sigma}}\,\xi^{s*}
  \end{pmatrix},
\end{equation}
其中我们用了 $v^s(p) = (u^s(p))^{c}$,即电荷共轭。也就是说,
\begin{equation}
  u^s(p) = -i\gamma^2(v^s(p))^{*},\qquad
  v^s(p) = -i\gamma^2(u^s(p))^{*}.
  \label{eq:3.135}
\end{equation}
如果我们把 \eqref{eq:3.135} 代入费米子场算符的表达式,然后用 $C$ 变换这个算符,我们得到
\begin{equation}
  C\psi(x)C = -i\gamma^2\psi^{*}(x) = -i\gamma^2(\psi^{\dagger}\gamma^0)^T
    = -i(\bar{\psi}\gamma^0\gamma^2)^T .
  \label{eq:3.136}
\end{equation}
注意 $C$ 是一个线性幺正算符,尽管它把 $\psi$ 变为 $\psi^{c}$。

我们再一次想要知道 $C$ 如何作用于费米子双线性量。首先我们需要
\begin{equation}
  C\bar{\psi}(x)C = C\psi^{\dagger}(x)C\gamma^0
    = (-i\gamma^2\psi^{*})^{\dagger}\gamma^0
    = (-i\gamma^0\gamma^2\psi)^{\dagger} .
  \label{eq:3.137}
\end{equation}
算出双线性量的变换有点棘手,写出旋量指标会有帮助。对标量,
\begin{align}
  C\bar{\psi}\psi\,C &= (-i\bar{\psi}\gamma^0\gamma^2)^T(-i\gamma^2\psi^{*})
    = -\psi^{\dagger}\gamma^2\gamma^0\gamma^2\psi^{*}
    = \psi^{\dagger}\gamma^0\psi
    = +\bar{\psi}\psi .
  \label{eq:3.138}
\end{align}
(第三步中的负号来自费米子反对易。)赝标量并不更难:
\begin{equation}
  Ci\bar{\psi}\gamma^5\psi\,C = +i\bar{\psi}\gamma^5\psi .
  \label{eq:3.139}
\end{equation}
我们必须分别处理矢量和赝矢量的每个分量。注意到 $\gamma^0$ 和 $\gamma^2$ 是对称矩阵,而 $\gamma^1$ 和 $\gamma^3$ 是反对称矩阵,我们最终得到
\begin{equation}
  C\bar{\psi}\gamma^\mu\psi\,C = -\bar{\psi}\gamma^\mu\psi;
  \label{eq:3.140}
\end{equation}
\begin{equation}
  C\bar{\psi}\gamma^\mu\gamma^5\psi\,C = +\bar{\psi}\gamma^\mu\gamma^5\psi .
  \label{eq:3.141}
\end{equation}
虽然算符 $C$ 互换 $\psi$ 和 $\bar{\psi}$,但它实际上并不改变产生和湮灭算符的顺序。因此,如果 $\bar{\psi}\gamma^0\psi$ 的定义中减去了式 \eqref{eq:3.113} 上方提到的无穷大常数,那么这个常数在 $C$ 共轭的过程中不会重新出现。

\subsection{C、P、T 总结}

各种费米子双线性量在 $C$, $P$ 和 $T$ 下的变换性质总结在下表中。这里我们用简写:对 $\mu = 0$, $(-1)^\mu = 1$;对 $\mu = 1, 2, 3$, $(-1)^\mu = -1$。
\begin{equation}
  \begin{array}{c|ccccc}
    & \bar{\psi}\psi & \bar{\psi}\gamma^5\psi
      & \bar{\psi}\gamma^\mu\psi & \bar{\psi}\gamma^\mu\gamma^5\psi
      & \bar{\psi}\sigma^{\mu\nu}\psi\\
    \hline
    P & +1 & -1 & (-1)^\mu & -(-1)^\mu & (-1)^\mu(-1)^\nu\\
    T & +1 & -1 & (-1)^\mu & (-1)^\mu & -(-1)^\mu(-1)^\nu\\
    C & +1 & +1 & -1 & +1 & -1\\
    CPT & +1 & +1 & -1 & -1 & +1
  \end{array}
  \label{eq:3.142}
\end{equation}
我们还包括了张量双线性量(见习题 3.7)以及导数算符的变换性质。

首先注意,自由狄拉克拉格朗日量 $\mathcal{L}_0 = \bar{\psi}(i\slashed{\partial} - m)\psi$ 在 $C$, $P$ 和 $T$ 下分别不变。我们可以在 $\mathcal{L}_0$ 中加入某个微扰 $\delta\mathcal{L}$,来构造破坏这些对称性中任意一个的更一般的量子系统。但 $\delta\mathcal{L}$ 必须是洛伦兹标量,而表的最后一行表明,$\psi$ 和 $\bar{\psi}$ 的所有洛伦兹标量组合在组合对称性 $CPT$ 下都不变。实际上,一般来说,人们无法构造一个具有厄米哈密顿量但破坏 $CPT$ 的洛伦兹不变的量子场论。%
\footnote{这一定理与自旋-统计定理在 Streater 和 Wightman, \emph{同上引文} 中有非常仔细的证明。}

\problems

\textbf{3.1} \emph{洛伦兹群}.  回忆式~\eqref{eq:3.17} 的洛伦兹对易关系,
\begin{equation}
  [J^{\mu\nu}, J^{\rho\sigma}] = i\bigl( g^{\nu\rho}J^{\mu\sigma}
    - g^{\mu\rho}J^{\nu\sigma} - g^{\nu\sigma}J^{\mu\rho}
    + g^{\mu\sigma}J^{\nu\rho} \bigr).
\end{equation}
\begin{enumerate}[label=(\alph*)]
  \item 定义旋转和洛伦兹加速的生成元为
    \begin{equation}
      J^i = \tfrac{1}{2}\epsilon^{ijk}J^{jk},\qquad
      K^i = J^{0i},
    \end{equation}
    其中 $i, j, k = 1, 2, 3$。无穷小洛伦兹变换于是可以写为
    \begin{equation}
      \Phi \to (1 - i\bm{\theta}\cdot\mathbf{J} - i\bm{\beta}\cdot\mathbf{K})\Phi .
    \end{equation}
    显式写出这些矢量算符的对易关系。(例如,$[J^i, J^j] = i\epsilon^{ijk}J^k$, $[K^i, J^j] =
    i\epsilon^{ijk}K^k$, $[K^i, K^j] = -i\epsilon^{ijk}J^k$。) 证明组合
    \begin{equation}
      \mathbf{J}^{+} = \tfrac{1}{2}(\mathbf{J} + i\mathbf{K})
      \qquad\text{以及}\qquad
      \mathbf{J}^{-} = \tfrac{1}{2}(\mathbf{J} - i\mathbf{K})
    \end{equation}
    彼此对易,并分别满足角动量的对易关系。
  \item 旋转群的有限维表示恰好对应角动量的允许值:整数或半整数。(a) 部分的结果意味着洛伦兹群的所有有限维表示对应整数或半整数的对 $(j_+, j_-)$,即旋转群的表示对。利用在角动量的自旋 $1/2$ 表示中 $J = \sigma/2$ 这一事实,显式写出按洛伦兹群 $(\tfrac{1}{2}, 0)$ 和 $(0, \tfrac{1}{2})$ 表示变换的 2 分量对象的变换律。证明这些恰好对应 \eqref{eq:3.39} 中给出的 $\psi_L$ 和 $\psi_R$ 的变换。
  \item 恒等式 $\sigma^T = -\sigma^2\sigma\sigma^2$ 使我们能够把 $\psi_L$ 变换改写成幺正等价的形式
    \begin{equation}
      \psi_L \to \psi_L'\Bigl( 1 + i\bm{\theta}\cdot\tfrac{\bm{\sigma}}{2}
        + \bm{\beta}\cdot\tfrac{\bm{\sigma}}{2} \Bigr),
    \end{equation}
    其中 $\psi_L' = \psi_L^{\dagger}\sigma^2$。利用这一变换律,我们可以把按 $(\tfrac{1}{2}, \tfrac{1}{2})$ 变换的对象表示为一个 $2\times2$ 矩阵,其左边具有 $\psi_R$ 的变换律,同时右边具有转置的 $\psi_L$ 变换。把这个矩阵参数化为
    \begin{equation}
      V = \begin{pmatrix}
        V^0 + V^3 & V^1 - iV^2\\
        V^1 + iV^2 & V^0 - V^3
      \end{pmatrix}.
    \end{equation}
    证明对象 $V^\mu$ 按 4-矢量变换。
\end{enumerate}

\textbf{3.2} 推导戈登恒等式,
\begin{equation}
  \bar{u}(p')\gamma^\mu u(p) = \bar{u}(p')\Bigl(
    \frac{p'^{\mu} + p^{\mu}}{2m} + \frac{i\sigma^{\mu\nu}q_\nu}{2m}
    \Bigr)u(p),
\end{equation}
其中 $q = (p' - p)$。我们将在第 6 章用到这一公式。

\textbf{3.3} \emph{旋量乘积}.  (本题连同习题 5.3 和 5.6 一起,引入了一种计算涉及无质量粒子过程的有效方法。) 设 $k_0$, $k_1$ 为满足 $k_0^2 = 0$, $k_1^2 = -1$, $k_0\cdot k_1 = 0$ 的固定 4-矢量。按如下方式定义基本旋量:设 $u_{L0}$ 为动量 $k_0$ 的费米子的左手旋量。设 $u_{R0} = k_{1L}u_{L0}$。然后,对任意类光($p^2 = 0$)的 $p$,定义
\begin{equation}
  u_L(p) = \frac{1}{\sqrt{2p\cdot k_0}}\,p_{R}u_{L0},\qquad
  \text{以及}\qquad
  u_R(p) = \frac{1}{\sqrt{2p\cdot k_0}}\,p_{L}u_{R0}.
\end{equation}
这一组约定无歧义地定义了旋量的相位(除非 $p$ 平行于 $k_0$)。
\begin{enumerate}[label=(\alph*)]
  \item 证明 $\slashed{k}_0 u_{R0} = 0$。证明对任意类光 $p$,有 $\slashed{p}u_L(p) = \slashed{p}u_R(p) = 0$。
  \item 对选择 $k_0 = (E, 0, 0, -E)$, $k_1 = (0, 1, 0, 0)$,显式构造 $u_{R0}$, $u_L(p)$ 和 $u_R(p)$。
  \item 对类光的 $p_1, p_2$,定义旋量乘积 $s(p_1, p_2)$ 和 $t(p_1, p_2)$ 为
    \begin{equation}
      s(p_1, p_2) = \bar{u}_R(p_1)u_L(p_2),\qquad
      t(p_1, p_2) = \bar{u}_L(p_1)u_R(p_2) .
    \end{equation}
    利用 (b) 部分给出的 $u$ 的显式形式,显式计算旋量乘积,并证明
    $t(p_1, p_2) = (s(p_2, p_1))^{*}$ 和 $s(p_1, p_2) = -s(p_2, p_1)$。
    此外,证明
    \begin{equation}
      |s(p_1, p_2)|^2 = 2p_1\cdot p_2 .
    \end{equation}
    因此旋量乘积是 4-矢量点积的平方根。
\end{enumerate}

\textbf{3.4} \emph{马约拉纳费米子}.  回忆式~\eqref{eq:3.46},可以对一个按狄拉克旋量 $(\tfrac{1}{2}, 0)$ 的上两个分量变换的无质量 2 分量费米子场写出相对论性方程。称这样的 2 分量场为 $\chi_\alpha(x)$,
$\alpha = 1, 2$。
\begin{enumerate}[label=(\alph*)]
  \item 证明可以按如下方式为 $\chi(x)$ 写出一个作为有质量场的方程:
    \begin{equation}
      i\bar{\sigma}\cdot\partial\,\chi - im\sigma^2\chi^{*} = 0 .
    \end{equation}
    也就是说,首先证明该方程是相对论不变的,其次证明它蕴含克莱因-戈登方程,
    $(\partial^2 + m^2)\chi = 0$。这种形式的费米子质量称为
    \emph{马约拉纳质量项}。
  \item 马约拉纳方程能否由拉格朗日量推出?质量项看起来是 $(\sigma^2)_{ab}\chi_a\chi_b$ 的变分;然而,
    由于 $\sigma^2$ 是反对称的,如果 $\chi(x)$ 是普通的 c-数场,这个表达式将为零。当我们进入量子场论时,我们知道 $\chi(x)$ 将变成一个反对易的量子场。因此,把 $\chi(x)$ 看作经典的 \emph{反对易}场来发展其经典理论是有意义的,也就是说,把它看作取值于满足下列关系的格拉斯曼数的场
    \begin{equation}
      \alpha\beta = -\beta\alpha
    \end{equation}
    对任意 $\alpha, \beta$。注意这个关系意味着
    $\alpha^2 = 0$。格拉斯曼场 $\zeta(x)$ 可以按一组函数基展开为
    \begin{equation}
      \zeta(x) = \sum_n a_n\phi_n(x),
    \end{equation}
    其中 $\phi_n(x)$ 是正交的 c-数函数,$a_n$ 是一组独立的格拉斯曼数。定义格拉斯曼数乘积的复共轭颠倒顺序:
    \begin{equation}
      (\alpha\beta)^{*} = \beta^{*}\alpha^{*} = -\alpha^{*}\beta^{*} .
    \end{equation}
    这条规则模仿了量子场的厄米共轭。证明经典作用量,
    \begin{equation}
      S = \int d^4x\ \Bigl[ \chi^{\dagger}i\bar{\sigma}\cdot\partial\chi
        + \tfrac{im}{2}\bigl( \chi^T\sigma^2\chi - \chi^{\dagger}\sigma^2
          \chi^{*} \bigr) \Bigr]
    \end{equation}
    (其中 $\chi^{\dagger} = (\chi^{*})^T$)是实的($S^{*} = S$),并且对 $\chi$ 和 $\chi^{*}$ 变分此 $S$ 得到马约拉纳方程。
  \item 让我们把一个 4 分量狄拉克场写为
    \begin{equation}
      \psi(x) = \begin{pmatrix}\chi_1(x)\\ \chi_2(x)\end{pmatrix},
    \end{equation}
    并回忆 $\psi$ 的下分量按与表示 $(\tfrac{1}{2}, 0)$ 的复共轭经幺正变换等价的方式变换。这样,我们可以用两个 2 分量旋量改写 4 分量狄拉克场:
    \begin{equation}
      \psi_L(x) = \chi_1(x),\qquad
      \psi_R(x) = i\sigma^2\chi_2^{*}(x) .
    \end{equation}
    用 $\chi_1$ 和 $\chi_2$ 改写狄拉克拉格朗日量,并注意质量项的形式。
  \item 证明 (c) 部分的作用量具有整体对称性。计算流
    \begin{equation}
      J^\mu = \chi^{\dagger}\bar{\sigma}^\mu\chi,\qquad
      J^\mu = \chi_1^{\dagger}\bar{\sigma}^\mu\chi_1
        - \chi_2^{\dagger}\bar{\sigma}^\mu\chi_2,
    \end{equation}
    分别对 (b) 和 (c) 部分的理论,并将你的结果与这些理论的对称性联系起来。构造一个具有 $O(N)$ 对称性(即 $N$ 维空间中的旋转对称性)的 $N$ 个自由有质量 2 分量费米子场的理论。
  \item 量子化 (a) 和 (b) 部分的马约拉纳理论。也就是说,把 $\chi(x)$ 提升为满足正则反对易关系的量子场
    \begin{equation}
      \{\chi_\alpha(\mathbf{x}), \chi_\beta^{\dagger}(\mathbf{y})\}
        = \delta_{\alpha\beta}\delta^{(3)}(\mathbf{x}-\mathbf{y}),
    \end{equation}
    构造厄米哈密顿量,并找到用一组产生和湮灭算符对哈密顿量对角化的正则对易关系的表示。( \emph{提示}:将 $\chi(x)$ 与量子化狄拉克场的上两个分量作比较。)
\end{enumerate}

\textbf{3.5} \emph{超对称}.  可以写出具有联系费米子和玻色子的连续对称性的场论;这样的变换称为超对称。
\begin{enumerate}[label=(\alph*)]
  \item 超对称场论最简单的例子是一个自由复玻色子和一个自由外尔费米子的理论,写为
    \begin{equation}
      \mathcal{L} = \partial_\mu\phi^{*}\partial^\mu\phi
        + \chi^{\dagger}i\bar{\sigma}\cdot\partial\chi + F^{*}F .
    \end{equation}
    这里 $F$ 是一个辅助复标量场,其场方程为 $F = 0$。证明这个拉格朗日量在无穷小变换下(相差一个全散度)不变
    \begin{align}
      \delta\phi &= -i\epsilon^T\sigma^2\chi,\\
      \delta\chi &= \epsilon F + \sigma\cdot\partial\phi\,\sigma^2\epsilon^{*},\\
      \delta F &= -i\epsilon^{\dagger}\bar{\sigma}\cdot\partial\chi,
    \end{align}
    其中参数 $\epsilon_\alpha$ 是格拉斯曼数的 2 分量旋量。
  \item 证明项
    \begin{equation}
      \Delta\mathcal{L} = \bigl[ m\phi F + \tfrac{i}{2}m\chi^T\sigma^2\chi \bigr]
        + (\text{复共轭})
    \end{equation}
    也被 (a) 部分给出的变换保持不变。通过求解其场方程从完整的拉格朗日量 $\mathcal{L} +
    \Delta\mathcal{L}$ 中消去 $F$,并证明费米子场和玻色子场 $\phi$ 和 $\chi$ 被赋予相同的质量。
  \item 通过向拉格朗日量添加三次及更高阶的项,可以写出超对称的非线性场方程。证明下面这个相当一般的、包含场
    $(\phi_i, \chi_i)$, $i = 1, \ldots, n$ 的场论是超对称的:
    \begin{equation}
      \mathcal{L} = \partial_\mu\phi_i^{*}\partial^\mu\phi_i
        + \chi_i^{\dagger}i\bar{\sigma}\cdot\partial\chi_i + F_i^{*}F_i
        + \Bigl[ W_i F_i - \tfrac{1}{2}W_{ij}\chi_i\chi_j + \text{c.c.}
          \Bigr],
    \end{equation}
    其中 $W[\phi]$ 是 $\phi_i$ 的任意函数,称为
    \emph{超势},下标表示对 $\phi_i$ 求导。对简单情形 $n = 1$ 和 $W = g\phi^3/3$,写出 $\phi$ 和 $\chi$ 的场方程(在消去 $F$ 之后)。
\end{enumerate}

\textbf{3.6} \emph{费尔兹变换}.  设 $u_i$, $i = 1, \ldots, 4$, 为四个 4 分量狄拉克旋量。在正文中,我们证明了费尔兹重排公式 \eqref{eq:3.83} 和 \eqref{eq:3.84}。这些公式中的第一个可以用 4 分量记号写为
\begin{equation}
  (\bar{u}_1\gamma^\mu P_L u_2)(\bar{u}_3\gamma_\mu P_L u_4)
    = -(\bar{u}_1\gamma^\mu P_L u_4)(\bar{u}_3\gamma_\mu P_L u_2)
      + \cdots .
\end{equation}
事实上,对任何乘积都有类似的重排公式
\begin{equation}
  (\bar{u}_1\Gamma^A u_2)(\bar{u}_3\Gamma^B u_4),
\end{equation}
其中 $\Gamma^A, \Gamma^B$ 是 3.4 节中列出的狄拉克矩阵的 16 种组合中的任意两种。
\begin{enumerate}[label=(\alph*)]
  \item 首先,把 16 个矩阵 $\Gamma^A$ 归一化到约定
    \begin{equation}
      \mathrm{tr}[\Gamma^A\Gamma^B] = 4\delta^{AB} .
    \end{equation}
    这给出 $\Gamma^A = \{1, \gamma^0, i\gamma^1, \ldots\}$;写出这个集合的全部 16 个元素。
  \item 把一般的费尔兹恒等式写为方程
    \begin{equation}
      (\bar{u}_1\Gamma^A u_2)(\bar{u}_3\Gamma^B u_4) =
        \sum_{C,D} C^{ABCD}(\bar{u}_1\Gamma^C u_4)(\bar{u}_3\Gamma^D u_2),
    \end{equation}
    其中 $C^{ABCD}$ 是未知系数。利用 16 个 $\Gamma^A$ 矩阵的完备性,证明
    \begin{equation}
      C^{ABCD} = \tfrac{1}{16}\mathrm{tr}[\Gamma^C\Gamma^A\Gamma^D\Gamma^B] .
    \end{equation}
  \item 显式算出乘积 $(\bar{u}_1u_2)(\bar{u}_3u_4)$ 和
    $(\bar{u}_1\gamma^\mu u_2)(\bar{u}_3\gamma_\mu u_4)$ 的费尔兹变换律。
\end{enumerate}

\textbf{3.7} 本题涉及分立对称性 $P$, $C$ 和 $T$。
\begin{enumerate}[label=(\alph*)]
  \item 计算反对称张量费米子双线性量在 $P$, $C$ 和 $T$ 下的变换性质,
    \begin{equation}
      \bar{\psi}\sigma^{\mu\nu}\psi,\qquad
      \text{其中 } \sigma^{\mu\nu} = \tfrac{i}{2}[\gamma^\mu, \gamma^\nu] .
    \end{equation}
    这补全了章末双线性量变换性质的表。
  \item 设 $\phi(x)$ 是一个复值克莱因-戈登场,正如我们在习题 2.2 中考虑的那样。找出幺正算符 $P$, $C$ 和一个反幺正算符 $T$(全部通过它们对克莱因-戈登粒子和反粒子的湮灭算符 $a_{\mathbf{p}}$ 和 $b_{\mathbf{p}}$ 的作用来定义),它们给出克莱因-戈登场的下列变换:
    \begin{equation}
      P\phi(t,\mathbf{x})P = \phi(t,-\mathbf{x});\qquad
      T\phi(t,\mathbf{x})T = \phi(-t,\mathbf{x});\qquad
      C\phi(t,\mathbf{x})C = \phi^{*}(t,\mathbf{x}) .
    \end{equation}
    求流的分量
    \begin{equation}
      j^\mu = i\bigl( \phi^{*}\partial^\mu\phi - \phi\partial^\mu\phi^{*} \bigr)
    \end{equation}
    在 $P$, $C$ 和 $T$ 下的变换性质。
  \item 证明任何由 $\psi(x)$, $\phi(x)$ 及其共轭构成的厄米洛伦兹标量局域算符都有 $CPT = +1$。
\end{enumerate}

\textbf{3.8} \emph{束缚态}.  两个自旋 $1/2$ 粒子可以组合成总自旋为 0 或 1 的态。这两个态的波函数在两个自旋交换下分别是奇的和偶的。
\begin{enumerate}[label=(\alph*)]
  \item 利用这一信息计算具有 $S$, $P$ 或 $D$ 波函数的所有电子-正电子束缚态在 $P$ 和 $C$ 下的量子数。
  \item 由于电子-光子耦合由哈密顿量给出
    \begin{equation}
      H_I = -e\int d^3x\ \bar{\psi}\gamma^\mu\psi\,A_\mu,
    \end{equation}
    其中 $j^\mu$ 是电流,如果矢量势的分量与 $j^\mu$ 的对应分量具有相同的 $P$
    和 $C$ 宇称,那么电动力学对 $P$ 和 $C$ 不变。证明这意味着下面这个令人惊讶的事实:电子偶素（正电子素）的自旋 0 基态可以衰变为 2 个光子,但自旋 1 基态必须衰变为 3 个光子。求更高电子偶素（正电子素）态湮灭以及电子偶素（正电子素）能级之间单光子跃迁的选择定则。
\end{enumerate}
